\documentclass[12pt]{memoir}

% Metadata follows the English notes and worksheets in this notebook.
\def\nsemestre{II}
\def\nterm{September}
\def\nyear{2026}
\def\nprofesor{}
\def\nauthor{Ignacio Rojas}
\def\nsigla{BG}
\def\nsiglahead{The classifying stack}
\def\nextra{Expository notes}
\def\nlang{ENG}
\let\footruleskip\relax
\input{../../../headerVarillyDiff.tex}
\author{Ignacio Rojas\\[0.4ex]ChatGPT (OpenAI)}

% Only document-specific helpers: notation, colors, theorem environments,
% multicol, hyperref, and the TikZ cd library already come from the header.
\newcommand{\BGtag}[1]{\href{https://stacks.math.columbia.edu/tag/#1}{Tag~\texttt{#1}}}
\newcommand{\BGcheckpoint}[1]{\par\smallskip{\raggedright\noindent\blu{\textbf{Checkpoint.}} #1\par}\smallskip}
\newcommand{\BGsolution}[1]{\subsection*{Solution to Exercise~\ref{#1}}\label{sol:#1}}
\counterwithout{section}{chapter}
\setsecnumdepth{subsection}
\settocdepth{section}
\setlength{\columnsep}{20pt}
\setlength{\emergencystretch}{1.5em}
\renewcommand{\bibsection}{\section{References}\label{sec:references}}
\hypersetup{pdftitle={The classifying stack BG: from categories to descent},
  pdfauthor={Ignacio Rojas; ChatGPT (OpenAI)},
  pdfsubject={Fppf torsors, fibered categories, and effective descent}}

\begin{document}
\maketitle
\begin{center}
\emph{From an ordinary category to a stack, with double covers as a running example}
\par\bigskip
\begin{minipage}{0.8\textwidth}
\small\centering
Developed from Ignacio Rojas's mathematical questions and reading notes,
with exposition, drafting, and revision by ChatGPT (OpenAI).
The examples and proofs are intended to make the terminology usable
before it is compressed into general definitions.
\end{minipage}
\end{center}
\clearpage
\tableofcontents*
\clearpage
\begin{multicols}{2}
\raggedcolumns

\section{Introduction and purpose}\label{sec:intro}

The aim is to prove that $BG$ is a stack, starting with its objects and
arrows. We first display those data, so that every subsequent definition
has a specific purpose: explaining an ingredient of $(S,P)$ or an
operation on such pairs. The preliminary sections are an illustrated
guide to reading this definition, not a list of prerequisites to master
before meeting $BG$.

The reader is assumed to know basic schemes, ring maps, and categories.
Group schemes, torsors, sites, and the functor of points are developed
here through separate examples and diagrams. In particular, an abstract
two-element group will give a family of two labelled copies of each
base scheme. We begin the meaning of ``two copies'' with finite sets.

The exposition follows the categorical viewpoint in Fantechi's
\emph{Stacks for Everybody}, especially Sections~3--4 and~6.3
\cite{Fantechi}. Exact external results are cited from the Stacks Project
\cite{Stacks}. The final comparison explains how to transfer the argument
to families of curves. We prove the stack axioms, not the additional
property called algebraicity. Set-theoretic size conventions play no
role in the calculations: we use ``collection'' when the objects need
not form a set, just as when speaking of the category of all schemes.

\section{Starting point: the data defining BG}\label{sec:start}

Fix a scheme $B$ and a group scheme $G\to B$. The example throughout is
$G=(\bZ/2\bZ)_B$. The following definition names the data first; the
subsequent sections unpack \emph{$B$-scheme}, \emph{group scheme},
\emph{fppf torsor}, and \emph{equivariant}. All torsors in the definition
are schemes. The assumptions on $G$ used in the proof are specified
and explained in Section~\ref{sec:groups}.

\begin{Def}[Objects and arrows for $BG$]\label{def:BGdata}
The data from which we build the total category $BG$ are as follows.
\begin{itemize}
\item \emph{Objects:} pairs $(S,P)$, where $S\to B$ is a $B$-scheme
and $P\to S$ is an fppf torsor for $G_S=G\times_B S$. The notation
$P$ includes its structure map and its $G_S$-action.
\item \emph{Morphisms:} a pair $(f,\widetilde f)$ from $(S,P)$ to
$(T,Q)$ consists of a $B$-morphism $f:S\to T$ and an equivariant
scheme morphism $\widetilde f:P\to Q$ making the square commute:
\end{itemize}
\[
\begin{tikzcd}[column sep=large,row sep=large]
P \arrow[r,"\widetilde f"] \arrow[d] &Q \arrow[d]\\
S \arrow[r,"f"'] &T.
\end{tikzcd}
\]
\end{Def}

There are three levels in this definition:
\begin{itemize}
\item $B$ is the fixed base \emph{scheme}.
\item $\Sch/B$ is the base \emph{category}; its objects are $S\to B$.
\item $BG$ is the proposed \emph{total category}; its objects carry
both $S$ and a torsor $P\to S$.
\end{itemize}
Thus $(S,P)$ is not an object of the base category. The projection
will forget $P$ and retain $S$. We have specified objects and arrows;
identities, composition, and the category axioms will be checked in
Section~\ref{sec:total}, after the words in the definition have been
explained. Defining the data does not yet establish those axioms.

\subsection{Roadmap for the proof}
Here is the route we will follow to show that $BG$ is a stack.
\begin{enumerate}[label=(\arabic*)]
\raggedright
\item $BG$ is a category: define identities and composition and check
their laws.
\item $BG$ is a category over $\Sch/B$: verify the projection functor.
\item $BG$ is fibered over $\Sch/B$: construct pullbacks and prove
their cartesian universal property.
\item $BG$ is fibered in groupoids: prove that every arrow over an
identity in the base is invertible.
\item Isomorphisms between torsors form a sheaf: glue compatible local
isomorphisms uniquely.
\item Object descent is effective: construct a global torsor from
compatible local torsors.
\item Collect these results to conclude that $BG$ is a stack.
\end{enumerate}
Each numbered item needs its own verification; the last two gluing
steps are additional to constructing a fibered category.

\subsection{Guide to the ingredients}
\begin{itemize}
\item Section~\ref{sec:categories}: $B$-schemes, categories over a base,
disjoint unions, and fiber products.
\item Section~\ref{sec:points}: how to interpret points of a scheme as
maps, with affine coordinates and a proof of the Yoneda argument used.
\item Sections~\ref{sec:groups}--\ref{sec:torsors}: the group scheme
$G$, the covering rule, and the torsor $P\to S$.
\item Section~\ref{sec:quotient}: why the notation $[\operatorname{pt}/G]$
leads to these objects.
\end{itemize}
After this introduction to the ingredients, Section~\ref{sec:total}
returns to Definition~\ref{def:BGdata} and starts the proof in the
order above. Exercises~\ref{ex:coproducts}--\ref{ex:coveroperations}
provide additional practice with the preliminary constructions.

\section[Categories and categories over a base]{Categories over a base}\label{sec:categories}

\subsection{Schemes and schemes with a specified base}
A scheme $S$ has an underlying topological space and a structure sheaf,
locally isomorphic as a locally ringed space to $\Spec A$ for a ring $A$.
The topology and sheaf are not arbitrary extra choices. A
\emph{$B$-scheme} is this same kind of scheme together with a specified
morphism $s:S\to B$. The added structure is the map, not a new meaning
of the word scheme.

\begin{Ex}[The affine line as a scheme over a field]
Let $k$ be a field and take $B=\Spec k$. The inclusion of constants
$\iota:k\to k[t]$ induces a scheme morphism in the reverse direction:
\[
\begin{tikzcd}[column sep=large]
k \arrow[r,"\iota"] & k[t]
\end{tikzcd}
\qquad c\longmapsto c,
\]
\[
\begin{tikzcd}[column sep=large]
\bA^1_k=\Spec k[t] \arrow[r,"\pi=\Spec\iota"] & \Spec k.
\end{tikzcd}
\]
What does $\pi$ do? A topological point of $\Spec k[t]$ is a prime
ideal $\mathfrak p$. The underlying map sends it to
$\iota^{-1}(\mathfrak p)=\mathfrak p\cap k$. This is always $(0)$:
a proper ideal cannot contain a nonzero constant, which is a unit.
Since $(0)$ is the only prime ideal of a field, every point maps to
the unique point of $\Spec k$. For instance, $(t-a)$ maps to $(0)$.

The scheme morphism also includes a map on functions: a constant $c$
on the base pulls back to the constant polynomial $c$. Thus the map
of spectra and the inclusion of constants describe the same morphism
from its two sides. The general scheme--ring correspondence is recalled
in Section~\ref{sec:points}.
\end{Ex}

\begin{Ex}[The relative affine line and its local squares]
For a general base scheme $B$, the relative affine line comes with
a map $\pi:\bA^1_B\to B$. Over an affine open $U=\Spec A\subset B$,
its description is
\[
\begin{tikzcd}[column sep=small,row sep=large]
\Spec A[t] \arrow[r,hook] \arrow[d,"\pi_U"']
 &\bA^1_B \arrow[d,"\pi"]\\
U=\Spec A \arrow[r,hook] &B.
\end{tikzcd}
\]
The left vertical map is induced by $A\to A[t]$, $a\mapsto a$;
on prime ideals it sends $\mathfrak p$ to $\mathfrak p\cap A$.
Unlike the field example, the target can have many points.
The square says that $\Spec A[t]$ is the inverse image of $U$ under
$\pi$. It is a fiber-product square, in the sense explained below.

For a principal open $D(a)\subset U$, the smaller local square is
\[
\begin{tikzcd}[column sep=small,row sep=large]
\Spec A_a[t] \arrow[r,hook] \arrow[d]
 &\Spec A[t] \arrow[d]\\
\Spec A_a \arrow[r,hook] &\Spec A.
\end{tikzcd}
\]
The horizontal maps come from localization, and the vertical maps
from inclusion of constants. The identity
$A[t]_a=A_a[t]$ makes these descriptions agree on smaller opens.
They glue to the relative affine line. This is the affine-local
description of Definition~27.5.1, \BGtag{01M0}, with the restriction
and base-change properties in Lemma~27.4.6, \BGtag{01LX}.
\end{Ex}

\begin{Def}[The category of $B$-schemes]\label{def:slice}
The category $\Sch/B$ is defined by the following data.
\begin{itemize}
\item \emph{Objects:} pairs $(S,s)$ with $S$ a scheme and $s:S\to B$
a scheme morphism.
\item \emph{Morphisms:} an arrow $(S,s)\to(T,t)$ is a scheme morphism
$f:S\to T$ such that $t\circ f=s$.
\item \emph{Identities and composition:} use the identity and
composition of scheme morphisms. The equation defining an arrow is
preserved: if $u\circ g=t$ and $t\circ f=s$, then
$u\circ(g\circ f)=s$.
\end{itemize}
The condition on morphisms is the commuting triangle
\[
\begin{tikzcd}[column sep=small,row sep=large]
S \arrow[rr,"f"] \arrow[dr,"s"'] && T \arrow[dl,"t"]\\
& B &
\end{tikzcd}
\]
The identity and associativity laws are inherited from $\Sch$.
\end{Def}

\begin{Rmk}[Comparison with Fantechi]
This is exactly the slice construction in Fantechi's Section~3.2,
with $B$ in place of her fixed scheme $S$; it is
also called the \emph{slice category}. We shall often abbreviate its
object $(S,s)$ to $S$, but the map $s$ remains part of the data.
``A scheme over $B$'' and ``a $B$-scheme'' mean the same thing.
The object $(B,\id_B)$ is terminal: from $(S,s)$ there is exactly one
arrow to it, namely $s$.
\end{Rmk}

\subsection{Categories over a category, and their fibers}
\begin{Def}
A category consists of objects, sets of arrows between objects, identity
arrows, and an associative composition law with left and right identity
laws. A functor preserves the sources and targets of arrows, identities,
and composition.
\end{Def}

\begin{Def}
A \emph{category over} a category $\cC$ is a category $\cF$ equipped with a
functor $p:\cF\to\cC$. An object $x$ lies over $S$ if $p(x)=S$;
an arrow $u:x\to y$ lies over $f:S\to T$ if $p(u)=f$.
The \emph{fiber category} $\cF_S$ is specified by
\[
 \operatorname{Ob}(\cF_S)=\{x:p(x)=S\}
\]
and, for two such objects $x,y$, by
\[
\begin{aligned}
\Hom_{\cF_S}(x,y)
=\{u\in\Hom_{\cF}(x,y):\\[-0.4ex]
\hfill p(u)=\id_S\}.
\end{aligned}
\]
These arrows are called \emph{vertical}.
\end{Def}

We use $\cF_S$ for the fiber throughout, and $(BG)_S$ for the fiber
of $BG$. One may also write $p^{-1}(S)$ provided that its arrows mean
exactly those in the definition, not all arrows between the selected
objects. Later, parentheses in $X(T)$ will denote a \emph{set of points},
not a fiber category.

\begin{Rmk}[Vertical arrows]
``Vertical'' is a geometric nickname: imagine objects of $\cF$ arranged
above the objects of $\cC$. A vertical arrow does not move in the base;
its base arrow is the identity. It need not be drawn vertically on the
page. For a general functor there is no additional triangle to check:
$p(u)=\id_S$ is the complete condition. Triangles arise when the
particular categories have themselves been defined using triangles,
as with $\Sch/B$.
\end{Rmk}

\begin{Ej}[Fiber versus full subcategory]\label{ex:fiber}
A \emph{full subcategory} on selected objects retains every arrow of
the original category between those objects. Let $\cD$ be the full
subcategory of $\cF$ on the objects $x$ with $p(x)=S$.
Write the object collections and arrow sets of $\cD$ and $\cF_S$
separately. Then take $p=\id_{\cC}$ and an endomorphism
$\varepsilon:S\to S$ with $\varepsilon\ne\id_S$.
Which category contains $\varepsilon$? Why? Use $t\mapsto t^2$ on
$\bA^1_k$ as a concrete example. The full comparison is worked out in
the \hyperref[sol:ex:fiber]{solution on page~\pageref*{sol:ex:fiber}}.
\end{Ej}

\subsection{Fiber products: construction, then universal property}
Suppose given $q:X\to S$ and $f:T\to S$. We want a scheme over $T$
that retains precisely the part of $X$ lying over the base specified by
$f$. This operation is called \emph{base change} or \emph{pullback}.
\begin{Ex}[A fiber product of finite sets]\label{ex:finitefiber}
Take three sets
\[
\begin{gathered}
 S=\{0,1\},\qquad T=\{a,b,c\},\\
 X=\{u,v,w\}.
\end{gathered}
\]
Define maps by
\[
\begin{array}{c|ccc}
t&a&b&c\\\hline
f(t)&0&0&1
\end{array}
\qquad
\begin{array}{c|ccc}
x&u&v&w\\\hline
q(x)&0&1&1.
\end{array}
\]
The fiber product retains only pairs with the same image in $S$:
\[
\begin{aligned}
T\times_S X
 &=\{(t,x):f(t)=q(x)\}\\
 &=\{(a,u),(b,u),(c,v),(c,w)\}.
\end{aligned}
\]
For example, $(a,v)$ is excluded because $0\ne1$.
The projections remember the first and second entries:
\[
\begin{tikzcd}[column sep=large,row sep=large]
T\times_S X \arrow[r,"\pr_X"] \arrow[d,"\pr_T"']
 &X \arrow[d,"q"]\\
T \arrow[r,"f"'] &S.
\end{tikzcd}
\]
Over $a$ there is one choice, $u$; over $c$ there are two choices,
$v$ and $w$. Pullback replaces the old base labels $0,1$ by the new
base labels $a,b,c$, using $f$ to decide which old fiber to copy.
\end{Ex}

\begin{Ex}[A fiber product of affine schemes]\label{ex:affinefiber}
For schemes, structure sheaves must be retained as well. If
$S=\Spec R$, $T=\Spec A$, and $X=\Spec C$, use the two ring maps
$\alpha:R\to A$ and $\gamma:R\to C$. Form
\[
 T\times_S X=\Spec(A\otimes_R C).
\]
The ring diagram is
\[
\begin{tikzcd}[column sep=large,row sep=large]
R \arrow[r,"\gamma"] \arrow[d,"\alpha"']
 &C \arrow[d,"c\mapsto1\otimes c"]\\
A \arrow[r,"a\mapsto a\otimes1"'] &A\otimes_R C.
\end{tikzcd}
\]
Applying $\Spec$ reverses all four arrows:
\[
\begin{tikzcd}[column sep=small,row sep=large]
\Spec(A\otimes_R C) \arrow[r,"\pr_X"] \arrow[d,"\pr_T"']
 &\Spec C \arrow[d,"q"]\\
\Spec A \arrow[r,"f"'] &\Spec R.
\end{tikzcd}
\]
The ring square commutes because
$\alpha(r)\otimes1=1\otimes\gamma(r)$. The scheme square therefore
commutes too. For example, if $R=k$, $A=k[v]$, and $C=k[x]$, then
$A\otimes_R C\cong k[v,x]$: the product of two affine lines over
$k$ is the affine plane, with its coordinate projections.
Lemma~26.17.2, \BGtag{01JQ}, proves the general affine formula.
\end{Ex}

These affine constructions glue to give fiber products of arbitrary
schemes (Section~26.17, \BGtag{01JO}). If the original arrows are
$B$-morphisms, the product is a $B$-scheme via its map to $S\to B$.

\begin{Def}[The universal property of a fiber product]
The universal property describes exactly the maps \emph{into this
constructed scheme}. For every scheme $W$ with maps $b:W\to T$ and
$a:W\to X$ satisfying $fb=qa$, there exists exactly one map
$\lambda:W\to T\times_S X$ with
\[
 \pr_T\lambda=b,\qquad \pr_X\lambda=a.
\]
The diagram is
\[
\begin{tikzcd}[column sep=small,row sep=large]
W \arrow[drr,bend left,"a"] \arrow[ddr,bend right,"b"']
 \arrow[dr,dashed,"\lambda"] &&\\
&T\times_S X \arrow[r,"\pr_X"] \arrow[d,"\pr_T"']
 &X \arrow[d,"q"]\\
&T \arrow[r,"f"'] &S
\end{tikzcd}
\]
\end{Def}
Thus a compatible pair of maps from $W$ gives a \emph{unique map},
not generally an isomorphism.

\begin{Ex}[The unique map need not be an isomorphism]\label{ex:originplane}
Let the common target be $S=\Spec k$, and take
$T=\Spec k[v]$ and $X=\Spec k[x]$. Put $W=\Spec k$.
Define $b:W\to T$ by $v\mapsto0$ on rings and $a:W\to X$ by
$x\mapsto0$. Both composites $W\to S$ are the identity. The resulting
diagram is
\[
\begin{tikzcd}[column sep=small,row sep=large]
\Spec k \arrow[drr,bend left,"a"]
 \arrow[ddr,bend right,"b"'] \arrow[dr,dashed,"\lambda"] &&\\
&\bA^2_k \arrow[r,"\pr_X"] \arrow[d,"\pr_T"']
 &\bA^1_{x,k} \arrow[d]\\
&\bA^1_{v,k} \arrow[r] &\Spec k.
\end{tikzcd}
\]
The map $\lambda$ is the origin in the plane: its ring map is
$k[v,x]\to k$, $F\mapsto F(0,0)$. There is exactly one map with
these two coordinate maps, because the images of $v$ and $x$ determine
the ring homomorphism. It is not an isomorphism: that homomorphism has
the nonzero kernel $(v,x)$.

The point $W$ does not itself have the universal property. For instance,
the two maps from $Z=\Spec k[z]$ given by $v\mapsto z$ and
$x\mapsto0$ are compatible over $k$, but cannot factor through $W$
with its origin projections. Such a factorization would force the
image of $v$ to be $0$, not $z$.
\end{Ex}

\begin{Prop}[Uniqueness of a fiber product]\label{prop:fiberunique}
Suppose $F$, with projections $p_T,p_X$, and $W'$, with projections
$b',a'$, both satisfy the fiber-product universal property for
$T\xrightarrow{f}S\xleftarrow{q}X$. There is a unique isomorphism
$\alpha:W'\to F$ preserving both projections.
\end{Prop}
\begin{proof}
First use $W'$ as the source in the universal property of $F$.
It gives the dashed map $\alpha$ in
\[
\begin{tikzcd}[column sep=large,row sep=large]
W' \arrow[rr,"a'"] \arrow[dr,dashed,"\alpha"]
 \arrow[dd,"b'"'] &&X \arrow[dd,"q"]\\
&F \arrow[ur,"p_X"'] \arrow[dl,"p_T"] &\\
T \arrow[rr,"f"'] &&S.
\end{tikzcd}
\]
Thus $p_T\alpha=b'$ and $p_X\alpha=a'$.
Now use $F$ as the source in the universal property of $W'$:
\[
\begin{tikzcd}[column sep=large,row sep=large]
F \arrow[rr,"p_X"] \arrow[dr,dashed,"\beta"]
 \arrow[dd,"p_T"'] &&X \arrow[dd,"q"]\\
&W' \arrow[ur,"a'"'] \arrow[dl,"b'"] &\\
T \arrow[rr,"f"'] &&S.
\end{tikzcd}
\]
This gives $b'\beta=p_T$ and $a'\beta=p_X$.
Consequently
\[
 p_T\alpha\beta=p_T,\qquad p_X\alpha\beta=p_X.
\]
Both $\alpha\beta$ and $\id_F$ are maps $F\to F$ with these
projections. Uniqueness in the universal property of $F$ gives
$\alpha\beta=\id_F$. The same argument with $W'$ gives
$\beta\alpha=\id_{W'}$. Thus $\alpha$ and $\beta$ are inverse.
The first universal property already says that $\alpha$ is the only
map preserving the projections, hence the only such isomorphism.
\end{proof}

\begin{Rmk}
Merely having the outer commuting square for $W'$ gives $\alpha$,
but not $\beta$. The second universal property, valid for \emph{every}
source scheme, is what supplies the inverse.
\end{Rmk}

\begin{Ex}[Pulling back a parabola]\label{ex:parabola}
Let $k$ be a field. Name the original base and the family:
\[
\begin{aligned}
 S&=\Spec k[t],\\
 X&=\Spec\bigl(k[t,x]/(x^2-t)\bigr).
\end{aligned}
\]
Thus $X$ is the closed subscheme $V(x^2-t)$ of the affine $(t,x)$-plane.
The map $q:X\to S$ is induced by
$k[t]\to k[t,x]/(x^2-t)$, $t\mapsto\bar t$. On $k$-valued points
it is precisely $(t,x)\mapsto t$. Over a value $t_0$, its fiber has
equation $x^2=t_0$. For $k=\bR$ the picture is the parabola $t=x^2$
projecting to the $t$-axis.

Now choose a \emph{new base} $T=\Spec k[v]$ and a map $f:T\to S$
whose ring homomorphism is
\[
 \alpha:k[t]\longrightarrow k[v],\qquad t\longmapsto v^2.
\]
On $k$-valued points, $f$ sends $v$ to $v^2$. The pullback must assign
to $v$ the old fiber over $t=v^2$: its equation is therefore $x^2=v^2$.
Here is exactly how the earlier tensor product produces that equation.
In Example~\ref{ex:affinefiber}, take
\[
\begin{gathered}
 R=k[t],\qquad A=k[v],\\
 C=k[t,x]/(x^2-t).
\end{gathered}
\]
Use $\alpha$ above for $R\to A$, and $t\mapsto\bar t$ for
$R\to C$. Then
\[
\begin{aligned}
A\otimes_R C
 &\cong k[v,t,x]/(t-v^2,x^2-t)\\
 &\cong k[v,x]/(x^2-v^2).
\end{aligned}
\]
The two relations first identify the copies of $t$, then retain the
equation defining $X$. Eliminating $t$ gives the second isomorphism.
Write $X'=\Spec k[v,x]/(x^2-v^2)$. The scheme square is
\[
\begin{tikzcd}[column sep=large,row sep=large]
X' \arrow[r,"h"] \arrow[d,"q'"'] &X \arrow[d,"q"]\\
T \arrow[r,"f"'] &S.
\end{tikzcd}
\]
The maps, on $k$-valued points, are
\[
\begin{aligned}
h(v,x)&=(v^2,x),& q'(v,x)&=v,\\
q(t,x)&=t,& f(v)&=v^2.
\end{aligned}
\]
In particular, $qh(v,x)=v^2=fq'(v,x)$.
On coordinate rings, $h$ sends $\bar t$ to $v^2$ and $\bar x$ to $x$.

The real picture suggested by the equations is accurate: the original
family is a parabola, while the pulled-back family is the union of
the lines $x=v$ and $x=-v$. But the operation is not replacing one
curve by an isomorphic curve. It changes the base map and retains a
comparison map $h$ to the original family. For instance, over $v=2$
the new fiber copies the old fiber over $t=4$, with choices $x=2,-2$.

If $\operatorname{char}k\ne2$, restrict to $v\ne0$. In
$k[v,v^{-1},x]$, the ideals $(x-v)$ and $(x+v)$ are comaximal:
their difference contains the unit $2v$. The Chinese remainder theorem
therefore gives
\[
 \frac{k[v,v^{-1},x]}{(x^2-v^2)}
 \cong k[v,v^{-1}]\times k[v,v^{-1}].
\]
This is where the two lines become two \emph{disjoint} copies of the
new base. Over $v=0$ they meet, so the full family is not the trivial
double cover. In characteristic $2$, the equation is $(x-v)^2=0$;
there are not two distinct lines. These qualifications matter when
we later compare double covers with torsors.
\end{Ex}

\subsection{Two copies are a disjoint union, not a product}
\begin{Ex}[Two labelled copies of a three-element set]\label{ex:labelledset}
Let $E=\{1,2,3\}$. The disjoint union $E\amalg E$ is
\[
\begin{aligned}
E\amalg E&=(E\times\{0\})\cup(E\times\{1\})\\
 &=\{1_0,2_0,3_0,1_1,2_1,3_1\},
\end{aligned}
\]
where $n_i$ abbreviates $(n,i)$. In particular $1_0\ne1_1$: the label
is part of the element. This is different from $E\cup E=E$, which
has three elements, and from $E\times E$, which has nine.

To give a function $F:E\amalg E\to E$, specify its value on each
of the six labelled elements. Equivalently, give two functions
$f_0,f_1:E\to E$ and set
\[
 F(n_0)=f_0(n),\qquad F(n_1)=f_1(n).
\]
For example, choose $f_0=\id_E$ and $f_1(n)=4-n$. The function is
\[
\begin{array}{c|rrrrrr}
y&1_0&2_0&3_0&1_1&2_1&3_1\\\hline
F(y)&1&2&3&3&2&1.
\end{array}
\]
The particular choice $f_0=f_1=\id_E$ gives the \emph{fold map}
$\pi(n_i)=n$. It forgets only the label. Both $2_0$ and $2_1$, for
instance, map to $2$.
\end{Ex}

\begin{Def}[Disjoint union and its structure sheaf]
For a scheme $S$, construct $S\amalg S$ with underlying set
$|S|\times\{0,1\}$. Write its points as $s_i=(s,i)$, with $s\in|S|$.
An open subset is of the form $U_0\amalg U_1$, where each $U_i$ is
open in $S$. Define
\[
 \cO_{S\amalg S}(U_0\amalg U_1)
 =\cO_S(U_0)\times\cO_S(U_1),
\]
with restriction on each factor separately. The inclusions
$j_i:S\to S\amalg S$ send $s$ to $s_i$ and identify the sheaf on
the $i$th copy with $\cO_S$.
\end{Def}

\begin{Prop}[The constructed union is a scheme and a coproduct]
For a scheme $S$, the preceding ringed space is a scheme. Together
with $j_0,j_1$, it is a coproduct in $\Sch$: for every scheme $Y$,
restriction induces a bijection
\[
\begin{gathered}
\Hom(S\amalg S,Y)\\
\longrightarrow\Hom(S,Y)\times\Hom(S,Y),\\
F\longmapsto(Fj_0,Fj_1).
\end{gathered}
\]
\end{Prop}
\begin{proof}
The displayed sheaf is a sheaf because gluing and uniqueness hold on
each copy separately. Every point $s_i$ has an open neighbourhood
isomorphic, as a ringed space, to an affine open of $S$. Thus the
ringed space is locally an affine scheme, proving the first assertion.

Given morphisms $f_0,f_1:S\to Y$, define $F$ to equal $f_i$ on the
$i$th open copy. There is no overlap between the copies, hence no
compatibility condition. The maps on spaces and sheaves assemble
uniquely and are locally morphisms of schemes, so they define a unique
scheme morphism $F$. Conversely, any $F$ restricts to these two maps.
\end{proof}

The symbol $\amalg$ is used for the categorical coproduct. The
proposition checks that the labelled-copy construction realizes that
universal property for schemes; it is not a second, unrelated object.
The resulting map $F$ is often written $[f_0,f_1]$, called a
\emph{copairing}. Its domain is a union, not a product:
\[
\begin{tikzcd}[column sep=large,row sep=large]
S \arrow[r,"j_0"] \arrow[dr,"f_0"']
 &S\amalg S \arrow[d,dashed,"F"]
 &S \arrow[l,"j_1"'] \arrow[dl,"f_1"]\\
&Y&
\end{tikzcd}
\]
For $Y=S$ and $f_0=f_1=\id_S$, this constructs the scheme morphism
\[
 \pi=[\id_S,\id_S]:S\amalg S\longrightarrow S.
\]
Its underlying map is $s_i\mapsto s$, and its pullback of functions
on an open $U\subset S$ is $a\mapsto(a,a)$ in
$\cO_S(U)\times\cO_S(U)$. This specifies both parts of the scheme
morphism. Its defining equations are $\pi j_0=\pi j_1=\id_S$.
The notation $[\id_S,\id_S]$ does not mean a function taking values
in ordered pairs; the pair records the two restrictions of a single
function. Interchanging labels defines an automorphism
$\sigma(s_i)=s_{1-i}$, with $\sigma^2=\id$ and $\pi\sigma=\pi$.
The composite $S\amalg S\to S\to B$ makes the union a $B$-scheme,
and $\pi$ is an arrow both in $\Sch/B$ and in $\Sch/S$.

\begin{Ex}[Two affine lines, with their coordinate rings]\label{ex:twosheets}
Take $B=\Spec k$, $A=k[z]$, and $S=\Spec A=\bA^1_k$.
The disjoint union of two affine lines is
\[
 S\amalg S\cong\Spec(k[z]\times k[z]).
\]
Indeed the idempotents $(1,0)$ and $(0,1)$ split this spectrum into
disjoint open subsets. Localizing at either idempotent gives $k[z]$,
so these subsets have exactly the structure sheaves of the two affine
lines. This is the product-ring calculation in Lemma~10.21.2,
\BGtag{00ED}, including the localizations in its proof.

The fold and the two inclusions correspond to these ring maps:
\[
\begin{aligned}
\pi^*:k[z]&\longrightarrow k[z]\times k[z],\\
p(z)&\longmapsto(p(z),p(z)),\\
j_0^*:k[z]\times k[z]&\longrightarrow k[z],\\
(p(z),q(z))&\longmapsto p(z),\\
j_1^*:k[z]\times k[z]&\longrightarrow k[z],\\
(p(z),q(z))&\longmapsto q(z).
\end{aligned}
\]
For example, the two labelled points above $z=2$ have prime ideals
$(z-2)\times k[z]$ and $k[z]\times(z-2)$. Contracting either ideal
along $\pi^*$ gives $(z-2)$ in $k[z]$. The involution interchanges the
factors of the product ring.

Compare three operations on the same $S$:
\[
\begin{aligned}
S\amalg S&=\Spec(k[z]\times k[z]),\\
S\times_k S&=\Spec k[z_0,z_1],\\
S\times_S S&\cong\Spec k[z].
\end{aligned}
\]
They give two disjoint lines, a plane, and one line, respectively.
\end{Ex}

The objects in our running example have the form
\[
 (S,\pi:P\to S,\sigma).
\]
Here $S$ varies among $B$-schemes, $P$ is a scheme over $S$, and $\sigma$
is the sheet-exchanging involution. We will make ``double cover''
precise in Section~\ref{sec:torsors}: it means a cover that becomes
two copies after an allowed base change, not just a map that happens
to be two-to-one on some set of points.

The \emph{total} collection lets $S$ vary; for example it includes the
trivial cover of $\Spec k$ and the trivial cover of $\bA^1_k$ when
$B=\Spec k$. The collection over one fixed $S$ only includes covers
of that $S$. These are the object data of Definition~\ref{def:BGdata}
in the running example. With the arrows there specified and the
category axioms checked in Section~\ref{sec:total}, they give $BG$
and its fiber $(BG)_S$, respectively. The fiber keeps only arrows
over $\id_S$, not every arrow between those objects in the total category.

A \emph{nontrivial} double cover is one not isomorphic over $S$ to the
fold map. One example, justified in Section~\ref{sec:sites}, is
$\Spec\bC\to\Spec\bR$ with complex conjugation. This example also
warns against counting topological points: each spectrum has one point,
but the map is a double cover in the algebraic sense.

\BGcheckpoint{\begin{itemize}
\item Defined: $B$-schemes, $\Sch/B$, categories over a base, and fibers
$\cF_S$.
\item Constructed: labelled disjoint unions and affine fiber products.
\item Proved: two fiber products are uniquely isomorphic preserving
their projections.
\item Check: write the fold map on labelled elements and on functions.
\item Check: distinguish a map into a fiber product from an isomorphism
between two fiber products.
\end{itemize}}

\section{Test schemes and the functor of points}\label{sec:points}

\subsection{A test point is a morphism over the fixed base}
\begin{Def}[$T$-valued points]\label{def:testpoints}
Fix $B$-schemes $b_X:X\to B$ and $b_T:T\to B$. Define the set
\[
 X(T):=\Hom_B(T,X).
\]
An element $\xi\in X(T)$ is therefore a scheme morphism
$\xi:T\to X$ satisfying $b_X\circ\xi=b_T$. It is called a
\emph{$T$-valued point}, or a \emph{$T$-point}, of $X$.
\end{Def}
The condition ``over $B$'' is the commutative triangle
\[
\begin{tikzcd}[column sep=small,row sep=large]
T \arrow[rr,"\xi"] \arrow[dr,"b_T"'] && X \arrow[dl,"b_X"]\\
&B&
\end{tikzcd}
\]
We call the source $T$ a \emph{test scheme} when considering this set
of maps into $X$. For example, when $B=\Spec k$, choosing $T=\Spec k$
asks for rational points; choosing $T=\Spec k[\epsilon]/(\epsilon^2)$
allows maps that
also detect nilpotent functions, as Example~\ref{ex:dualpoints} computes.

\begin{Rmk}[Notation and the suppressed base]
The notation $X(T)$ is standard: it is exactly Poonen's
Definition~2.3.1 in \cite[Section~2.3]{Poonen}. The base is suppressed
in that notation, but remains part of both objects. Thus $X(T)$ means
maps over our specified $B$, not all scheme maps $T\to X$.
The longer notation $h_X(T)$ will denote the same set. The letter
$\xi$ will name a morphism; a letter such as $z$ will name a polynomial
variable. Keeping these roles separate helps in the computations below.
\end{Rmk}

\subsection{The functor of points: both assignments}
\begin{Def}[Functor of points]\label{def:pointfunctor}
For a fixed $B$-scheme $X$, its functor of points is
\[
 h_X:(\Sch/B)^{\mathrm{op}}\longrightarrow\mathsf{Sets}.
\]
On objects, it assigns $T\mapsto h_X(T)=X(T)$. On an arrow
$u:T'\to T$ of $B$-schemes, it assigns the function
\[
\begin{aligned}
 h_X(u):X(T)&\longrightarrow X(T'),\\
 \xi&\longmapsto\xi\circ u.
\end{aligned}
\]
Here $\mathsf{Sets}$ is the category of sets and functions, and
$\mathrm{op}$ records that arrows reverse direction.
\end{Def}
This is the functor in Stacks Project Example~4.3.4, \BGtag{001O},
applied to $\Sch/B$; see also Schemes, Section~26.13, \BGtag{01J5}.

Precomposition is pictured by
\[
\begin{tikzcd}[column sep=small,row sep=large]
T' \arrow[rr,"u"] \arrow[dr,"\xi\circ u"'] &&T \arrow[dl,"\xi"]\\
&X&
\end{tikzcd}
\]
All arrows are over $B$, so the composite is again an allowed test
point. The identity law is $\xi\circ\id_T=\xi$. If
$v:T''\to T'$, associativity gives
\[
 (\xi\circ u)\circ v=\xi\circ(u\circ v).
\]
Consequently
$h_X(v)\circ h_X(u)=h_X(u\circ v)$, which is the composition law
for this contravariant functor. We have specified the object assignment,
the arrow assignment, and checked both functor laws.

\subsection{Affine test schemes turn the definition into algebra}
\begin{Th}[Affine morphisms and ring maps]\label{th:affinepoints}
For commutative rings $A$ and $R$, taking the induced map on global
sections gives a bijection
\[
\begin{gathered}
 \Hom_{\Sch}(\Spec R,\Spec A)\\
 \xrightarrow{\sim}\Hom_{\mathrm{rings}}(A,R),\\
 \xi\longmapsto\xi^\sharp.
\end{gathered}
\]
The inverse sends a ring map $\theta:A\to R$ to the scheme morphism
$\Spec\theta:\Spec R\to\Spec A$.
\end{Th}
This is the affine case of Lemma~26.6.4, \BGtag{01I1};
Lemma~26.6.5, \BGtag{01I2}, packages it as an equivalence between
affine schemes and the opposite category of rings. This is a theorem
about schemes, not Yoneda's lemma. For orientation, the underlying
map sends a prime ideal $\mathfrak q\subset R$ to
$\theta^{-1}(\mathfrak q)\subset A$; localization of $\theta$ gives
the required maps on structure sheaves. The theorem asserts that this
procedure and passage to global sections undo each other.

Now specialize explicitly to $B=\Spec k$, where $k$ is a field,
$X=\Spec A$, and $T=\Spec R$, with specified $k$-algebra structures
on $A$ and $R$. Being a map \emph{over $\Spec k$} means that the
corresponding ring map respects $k$. The two triangles are
\[
\begin{tikzcd}[column sep=small,row sep=large]
\Spec R \arrow[rr,"\xi"] \arrow[dr] &&\Spec A \arrow[dl]\\
&\Spec k&
\end{tikzcd}
\]
and, with directions reversed,
\[
\begin{tikzcd}[column sep=small,row sep=large]
A \arrow[rr,"\xi^\sharp"] &&R\\
&k \arrow[ul] \arrow[ur]&
\end{tikzcd}
\]
Thus the theorem restricts to a bijection
\[
 X(\Spec R)\simeq\Hom_{k\text{-alg}}(A,R).
\]
The abbreviation $X(R)$ for $X(\Spec R)$ is also customary when the
base ring is understood. It does not mean that a ring has become a
scheme without applying $\Spec$.

\begin{Ex}[Recall: rational points]\label{ex:rationalpoints}
Over $B=\Spec k$, a \emph{$k$-rational point} of $X$ is a
$k$-morphism $\xi:\Spec k\to X$; their set is denoted $X(k)$.
For an affine scheme defined by polynomial equations, this recovers
solutions whose coordinates lie in $k$. The word ``rational'' here
does not require $k=\bQ$.

For instance, let $X=\bA^1_k=\Spec k[z]$. By
Theorem~\ref{th:affinepoints}, a $k$-point corresponds to a
$k$-algebra map $k[z]\to k$. Such a map is determined by one
element $a\in k$, the image of $z$, and is evaluation:
\[
 \xi_a^\sharp(F(z))=F(a).
\]
The single point of $\Spec k$ maps to the prime ideal $(z-a)$.
Therefore $\bA^1_k(k)\simeq k$. For a general affine
$X=\Spec k[z_1,\ldots,z_n]/I$, the same argument gives the tuples
$(a_1,\ldots,a_n)\in k^n$ annihilating every polynomial in $I$.
This is the coordinate interpretation recalled in
\cite[Section~2.3.1]{Poonen}.
\end{Ex}

\begin{Ex}[An affine line has points valued in any algebra]
Keep $B=\Spec k$ and $X=\Spec k[z]$, but now take
$T=\Spec R$ for any $k$-algebra $R$. Then
\[
 X(T)\simeq\Hom_{k\text{-alg}}(k[z],R)\simeq R.
\]
The second bijection sends $\xi^\sharp$ to $\xi^\sharp(z)$.
Conversely, $r\in R$ gives evaluation $F(z)\mapsto F(r)$, so
there is no condition on $r$. For example, with $R=k[t]$, the rule
$z\mapsto t^2$ determines a $T$-point $\xi:T\to X$.
This is a morphism from an entire affine line, not a choice of a
single topological point of $X$.
\end{Ex}

\begin{Ex}[Dual numbers detect what field-valued points miss]\label{ex:dualpoints}
Let $B=\Spec k$ and
$X=\Spec k[z]/(z^2)$. A test point with source $T=\Spec R$
corresponds to a $k$-algebra map
\[
 \xi^\sharp:k[z]/(z^2)\longrightarrow R.
\]
Its value $r=\xi^\sharp(\overline z)$ must satisfy $r^2=0$;
conversely, every such $r$ gives a map. Thus
\[
 X(\Spec R)\simeq\{r\in R:r^2=0\}.
\]
If $R=L$ is a field extension of $k$, the only possibility is $r=0$.
Field-valued points cannot distinguish $X$ from $\Spec k$ in this
example: each has exactly one $L$-point for every such $L$.

Instead use the \emph{dual numbers}
$R=k[\epsilon]/(\epsilon^2)$, where $\epsilon$ denotes the residue
class and is nonzero although $\epsilon^2=0$. An element is uniquely
$a+b\epsilon$ with $a,b\in k$. Its square vanishes precisely when
$a=0$. Therefore
\[
 X(\Spec R)\simeq\{b\epsilon:b\in k\}.
\]
In particular, the assignments $\overline z\mapsto0$ and
$\overline z\mapsto\epsilon$ give two different scheme morphisms
$\xi_0,\xi_1:T\to X$. Their underlying maps of topological spaces
are the same: both spaces have just one point. They differ on
functions, since $\xi_0^\sharp(\overline z)=0$ and
$\xi_1^\sharp(\overline z)=\epsilon$. Allowing nonreduced test
schemes detects this difference. It does not manufacture extra
topological points.
\end{Ex}

\subsection{Changing the target: naturality as a square}
Keep $X,Y$ over $B$, and let $f:X\to Y$ be a $B$-morphism.
For every test scheme $T$, postcomposition gives a function
\[
\begin{aligned}
 h_f(T):X(T)&\longrightarrow Y(T),\\
 \xi&\longmapsto f\circ\xi.
\end{aligned}
\]
This time the arrow $f:X\to Y$ and the induced function point in
the same direction; the reversal in $h_X$ concerned its source $T$.
For each $u:T'\to T$, there is a commutative square of \emph{sets}
\[
\begin{tikzcd}[column sep=large,row sep=large]
X(T) \arrow[r,"h_f(T)"] \arrow[d,"h_X(u)"']
 &Y(T) \arrow[d,"h_Y(u)"]\\
X(T') \arrow[r,"h_f(T')"'] &Y(T')
\end{tikzcd}
\]
Starting with $\xi\in X(T)$, the two routes give respectively
$(f\circ\xi)\circ u$ and $f\circ(\xi\circ u)$, which agree by
associativity. Equivalently,
\[
 h_f(T')(\xi\circ u)=h_f(T)(\xi)\circ u.
\]
A collection of functions $X(T)\to Y(T)$ is called a
\emph{natural transformation} when it commutes with precomposition
in this way for every $u$. Thus ``natural in $T$'' means this
explicit compatibility, not an unspecified property of the formulas.

\subsection{The Yoneda argument we will actually use}
\begin{Th}[Yoneda for the functors of points]\label{th:pointyoneda}
Suppose functions $\eta_T:X(T)\to Y(T)$ are given for every
$B$-scheme $T$, and for every $u:T'\to T$ satisfy
\[
 \eta_{T'}(\xi\circ u)=\eta_T(\xi)\circ u.
\]
There is a unique $B$-morphism $f:X\to Y$ such that
$\eta_T(\xi)=f\circ\xi$ for every $T$ and every $\xi\in X(T)$.
\end{Th}
This is Lemma~4.3.5, \BGtag{001P}, applied to the category $\Sch/B$.
Here the short proof is useful enough to give in full.
\begin{proof}
Use $X$ itself as the test scheme. Its identity morphism is an element
of $X(X)$. Define
\[
 f:=\eta_X(\id_X)\in Y(X)=\Hom_B(X,Y).
\]
This is already a genuine $B$-morphism by the definition of $Y(X)$.
We must show it gives all the prescribed functions, not just the one
at $X$. For any $\xi:T\to X$, naturality gives the square
\[
\begin{tikzcd}[column sep=large,row sep=large]
X(X) \arrow[r,"\eta_X"] \arrow[d,"h_X(\xi)"']
 &Y(X) \arrow[d,"h_Y(\xi)"]\\
X(T) \arrow[r,"\eta_T"'] &Y(T).
\end{tikzcd}
\]
Apply it to $\id_X$ in the upper left. The left-then-bottom route
gives $\eta_T(\id_X\circ\xi)=\eta_T(\xi)$; the
top-then-right route gives $f\circ\xi$. This proves existence.
If $f'$ induced the same functions, evaluation at $\id_X$ would
give $f'=f'\circ\id_X=\eta_X(\id_X)=f$. This proves uniqueness.
\end{proof}

\begin{Rmk}[How pointwise notation will be used]
In these notes the formula $(g,p)\mapsto g\cdot p$ means a rule
for maps from every $T$, compatible with precomposition. It does not
mean a rule only on the underlying topological points. Yoneda then
turns such a natural rule into a unique scheme morphism.
To check equality of two already defined scheme maps, it likewise
suffices to check every test scheme: using their source and its
identity point detects equality.

For fiber products the universal property gives
\[
 (T\times_S X)(W)\simeq T(W)\times_{S(W)}X(W).
\]
On the right is a fiber product of \emph{sets of morphisms}: its
elements are compatible pairs of maps from $W$. It is not the product
of the underlying topological spaces of the schemes.
\end{Rmk}

\BGcheckpoint{\begin{itemize}
\item Defined $X(T)=h_X(T)=\Hom_B(T,X)$, with its base triangle.
\item Defined $h_X$ on objects and arrows; verified its functor laws.
\item Computed rational, affine-line, and dual-number test points.
\item Expressed naturality by a commutative square and proved the
Yoneda recovery formula $f=\eta_X(\id_X)$.
\end{itemize}}

\section{Group schemes, built from examples}\label{sec:groups}

\subsection{An ordinary group is not automatically a scheme}
Let $H=\bZ/2\bZ=\{0,1\}$ be the abstract two-element group. This is a
set with an addition law. We do \emph{not} take a map from this bare
two-element set into $B$ and call it a group scheme. Instead we build
an actual scheme
\[
 H_B=B_0\amalg B_1\longrightarrow B,
\]
where each $B_a$ is a copy of $B$ and the map is the fold map of
Example~\ref{ex:twosheets}. Over every point of $B$ there is a copy
with label $0$ and a copy with label $1$. If $B$ has infinitely many
points, $H_B$ does too.

The labels tell us how to define a multiplication morphism. There are
four components in $H_B\times_B H_B$, indexed by $(a,b)\in H^2$.
Each is a copy of $B$. Map it by the identity on $B$ to the component
of $H_B$ labeled $a+b$ modulo $2$:
\[
\begin{array}{c|cccc}
(a,b)&(0,0)&(0,1)&(1,0)&(1,1)\\\hline
a+b&0&1&1&0.
\end{array}
\]
The unit map $e:B\to H_B$ is inclusion of $B_0$. Inversion is the
identity on both components because $-a=a$ in $H$. These are genuine
morphisms of schemes, defined component by component. This is our
first group scheme.

\subsection{What ``group object'' packages}
An ordinary group can be described by a multiplication
$H\times H\to H$, a unit map $\{*\}\to H$ selecting the identity,
and an inverse map $H\to H$. The group axioms are equations between
composites of these maps. To make a group object in $\Sch/B$, replace
sets by $B$-schemes, products by fiber products over $B$, and the
one-point terminal set by the terminal $B$-scheme $(B,\id_B)$.

\begin{Def}[Group scheme]\label{def:groupscheme}
A group scheme over $B$ is a $B$-scheme $\pi_G:G\to B$ equipped with
$B$-morphisms
\[
\begin{aligned}
 m&:G\times_B G\to G,\\
 e&:B\to G,\\
 \iota&:G\to G
\end{aligned}
\]
for multiplication, unit, and inverse. They satisfy associativity,
the two unit laws, and the two inverse laws.
\end{Def}
For precision, write $G^n$ for the $n$-fold product over $B$.
Associativity is the commutative square
\[
\begin{tikzcd}[column sep=large,row sep=large]
G^3 \arrow[r,"m\times\id_G"] \arrow[d,"\id_G\times m"']
 &G^2 \arrow[d,"m"]\\
G^2 \arrow[r,"m"'] &G
\end{tikzcd}
\]
The unit and inverse laws, with each pairing a map $G\to G^2$, are
\[
\begin{aligned}
m\langle e\pi_G,\id_G\rangle&=\id_G,\\
m\langle\id_G,e\pi_G\rangle&=\id_G,\\
m\langle\iota,\id_G\rangle&=e\pi_G,\\
m\langle\id_G,\iota\rangle&=e\pi_G.
\end{aligned}
\]
The phrase ``over $B$'' includes $\pi_Ge=\id_B$ and, for example,
$\pi_Gm=\pi_G\pr_1=\pi_G\pr_2$. Multiplication combines two elements
over the \emph{same} base; it does not combine unrelated fibers.
There need not be a group structure on the underlying topological set
of $G$.

For every $T\to B$, the set $G(T)=\Hom_B(T,G)$ is an ordinary group:
multiply $g,h:T\to G$ by $m\circ(g,h)$; its identity is
$T\to B\xrightarrow{e}G$; its inverse sends $g$ to $\iota g$.
The diagram equations give the usual group axioms on this set.
This is equivalent to Definition~39.4.1, \BGtag{022S}; the equivalence
uses precisely the Yoneda argument of Section~\ref{sec:points}.

For a finite abstract group $H$, the same component construction gives
$H_B=\coprod_{h\in H}B$ with products of labels given by the group
law. Its $T$-points are locally constant maps $T\to H$: the inverse
images of the labeled components are an open-and-closed partition of
$T$. If $T$ is nonempty and connected, this is just a choice of one
element of $H$. For disconnected $T$ there can be more choices.
This construction is called the \emph{constant group scheme}
(Example~39.5.6 in \BGtag{047F}).

\subsection{An infinite example: addition on the affine line}
Let $B=\Spec A$. Take $G=\Spec A[x]$, and define multiplication
by the ring map
\[
 A[x]\longrightarrow A[x_1,x_2],\qquad x\longmapsto x_1+x_2.
\]
The unit is given by $x\mapsto0$ in $A$; inversion by $x\mapsto-x$.
Their equations are the ordinary addition laws. This group scheme is
denoted $\mathbb G_{a,B}$. On $T=\Spec R$ it gives the additive group
$(R,+)$; in particular $\mathbb G_{a,k}(k)=(k,+)$, an infinite group
when $k$ is infinite. We have put a group law on a familiar scheme,
not replaced a group by an unrelated topological space.

Likewise $\mathbb G_{m,B}=\Spec A[x,x^{-1}]$ uses $x\mapsto x_1x_2$
for multiplication, $x\mapsto1$ for its unit, and $x\mapsto x^{-1}$
for inversion. Its $R$-points form the group $R^\times$. These
constructions glue over arbitrary $B$; see Examples~39.5.1 and~39.5.3
in \BGtag{047F}.

\subsection{Changing the base of a group scheme}
For a specified $s:S\to B$, form
\[
\begin{tikzcd}[column sep=large,row sep=large]
G_S=G\times_B S \arrow[r] \arrow[d,"\pi_{G_S}"']
 &G \arrow[d,"\pi_G"]\\
S \arrow[r,"s"'] &B
\end{tikzcd}
\]
The result is an $S$-scheme. It becomes a group scheme over $S$ by
pulling back the three structure maps. For example,
\[
 G_S\times_S G_S\simeq(G\times_B G)\times_B S,
\]
so $m_S$ is $m$ pulled back to $S$. On test points the formulas are
\[
\begin{aligned}
 m_S((g_1,t),(g_2,t))&=(m(g_1,g_2),t),\\
 e_S(t)&=(e(s(t)),t),\\
 \iota_S(g,t)&=(\iota(g),t).
\end{aligned}
\]
The shared $t$ is exactly the condition in the product over $S$.
The original equations pull back to the new equations, so these maps
give a group scheme (Lemma~39.4.2, \BGtag{022T}). Thus $G_S$ is not a
single ordinary group: it is a scheme over $S$ with a group law, and
$G_S(T)$ is an ordinary group for every $T\to S$.
For $H_B=B\amalg B$ this construction gives $H_S=S\amalg S$ with
the same label addition. For $\mathbb G_{a,B}$ it gives $\bA^1_S$
with addition.

\subsection{The geometric hypotheses}
These adjectives apply to the \emph{scheme morphism} $\pi_G:G\to B$,
not to an abstract set of group elements. Recall the following
definitions for any scheme morphism $f:X\to Y$.

\emph{Affine.} The inverse image of each affine open $\Spec A\subset Y$
is an affine scheme $\Spec C$. This is a relative property; $Y$ need
not itself be affine (Section~29.11, \BGtag{01S5}).

\emph{Locally of finite presentation.} Around every point of $X$ we
can choose affine opens $\Spec C\subset X$ and $\Spec A\subset Y$
with the first mapping into the second, such that
\[
 C\simeq A[x_1,\ldots,x_n]/(f_1,\ldots,f_r)
\]
as an $A$-algebra, for finite $n,r$. The coefficients belong to the
base ring $A$, not necessarily $\bC$. Thus the recollection ``finitely
many variables and a finitely generated ideal of relations'' is exactly
right. It is not a claim that the scheme has finitely many points or
that its group of sections is a finitely presented abstract group.
Finite presentation of a morphism without ``locally'' additionally
requires quasi-compactness and quasi-separatedness
(Definition~29.22.1, \BGtag{01TP}). These additional global conditions
are not prerequisites for the torsor proof; the local condition is the
one imposed on $G\to B$.

\emph{Flat.} A ring map $A\to C$ is flat if tensoring an exact sequence
of $A$-modules with $C$ preserves exactness. Since tensor product is
already right exact, this means it also preserves injections. A scheme
morphism is flat when its induced local-ring maps are flat; one can
equivalently check the corresponding ring maps on affine charts
(Section~29.26, \BGtag{01U2}). This controls how algebraic relations
behave in a family; it is not a synonym for smoothness. Free modules
are flat, which verifies all the examples needed immediately.

Check these conditions for $H_B$ with $H=\bZ/2\bZ$. Above $\Spec A$
its algebra is
\[
 A\times A\simeq A[z]/(z^2-z).
\]
Evaluation at $z=0,1$ gives the isomorphism. The two factors are
separated by the comaximal ideals $(z)$ and $(z-1)$, even in
characteristic $2$. It is affine, has one generator and one relation,
and is free of rank two as an $A$-module. Hence its map to $B$ is
affine, flat, and locally of finite presentation.
For $\mathbb G_{a,B}$, the algebra $A[x]$ has one generator and no
relations, and is free as an $A$-module with basis $1,x,x^2,\ldots$.
This example shows that finite presentation as an \emph{algebra} does
not mean finite generation as a \emph{module}.

\begin{Def}[Standing geometric assumptions]\label{def:standing}
From now on $G\to B$ is a group scheme whose structure morphism is
affine, flat, and locally of finite presentation.
\end{Def}
Its unit section makes it surjective: $\pi_Ge=\id_B$, so every base
point has the point supplied by $e$ above it. Affineness will let us
construct glued torsors using algebras. Flatness and local finite
presentation will make torsor projections allowable covering maps.
These are sufficient working hypotheses, not a claim of minimality.

\BGcheckpoint{\begin{itemize}
\item Constructed: $(\bZ/2\bZ)_B$, $\mathbb G_a$, and $\mathbb G_m$.
\item Defined: group-scheme structure maps and the three geometric
hypotheses.
\item Verified: the constant group satisfies these hypotheses.
\item Check: distinguish $H$, $H_B$, and $H_B(T)$.
\end{itemize}}

\section{Coverings, sites, and sheaves}\label{sec:sites}

\subsection{Start with ordinary open covers}
On a topological space, an open cover lets us work on smaller pieces
and compare answers on intersections. For example, continuous
functions on $U_i$ agreeing on $U_i\cap U_j$ give one continuous
function on $\bigcup_i U_i$. A site keeps this idea but allows the
maps used as local pieces to be more general than open inclusions.
It is not another scheme, variety, or topological space.

In these notes a \emph{site} means a category together with a rule
specifying which families of arrows $\{U_i\to U\}$ count as covers.
The covering-family formulation requires three rules:
\begin{enumerate}[label=(\arabic*)]
\item An isomorphism, by itself, is a cover.
\item Covers of covering pieces combine to give a cover of the
original object.
\item The fiber products needed to pull a cover back along any arrow
exist, and the resulting family covers the new base.
\end{enumerate}
This is Definition~7.6.2, \BGtag{00VH}. The covering rule is often
called a Grothendieck topology; formally these families specify a
basis for that topology. We need only the families and these rules.

Thus $(\Sch/B)_{\mathrm{fppf}}$ means precisely ``the category of
$B$-schemes, equipped with the fppf covering rule.'' Its objects and
arrows are still the objects and arrows of $\Sch/B$. In particular,
\emph{not every arrow must be fppf}; the extra rule only singles out
covering families. Nor have we changed the Zariski topology inside
each individual scheme. When this is called a \emph{big} site, the
word indicates that arbitrary $B$-schemes are admitted as objects,
not merely open subsets of $B$. No set-theoretic construction of a
``small model'' is needed to follow this proof.

\subsection{The fppf covering rule}
\begin{Def}\label{def:fppf}
An fppf cover of $S$ is a family of $B$-morphisms $\{f_i:S_i\to S\}$
such that each $f_i$ is flat and locally of finite presentation, and
their images together cover $S$.
\end{Def}
This is Definition~34.7.1, \BGtag{021M}. The letters abbreviate the
French words for faithfully flat and of finite presentation. For a
family the surjectivity requirement is \emph{joint}: an individual
$S_i$ need not cover all of $S$. A single surjective flat morphism
locally of finite presentation is a cover by itself.

An ordinary Zariski open covering is an fppf covering, but the allowed
maps need not be open immersions. For example,
\[
 U=\Spec\bC\longrightarrow S=\Spec\bR
\]
is an fppf cover: the algebra $\bC=\bR[z]/(z^2+1)$ is free of rank
two over $\bR$ and finitely presented, and the map is surjective.
It is not an isomorphism of schemes. Over the new base $U$, the cover
itself becomes
\[
 U\times_S U\simeq U\amalg U.
\]
Indeed, as a $\bC$-algebra using the first tensor factor,
\[
\begin{aligned}
\bC\otimes_{\bR}\bC
 &\simeq\bC[z]/(z^2+1)\\
 &\simeq\bC\times\bC,
\end{aligned}
\]
since $z^2+1=(z-i)(z+i)$ has distinct roots. Conjugating the second
factor exchanges the two components. This is a double cover that
splits after an allowed base change. It is not split over $\bR$:
a section would give an $\bR$-algebra map $\bC\to\bR$, impossible
because $-1$ has no square root in $\bR$. Also, $\Spec\bR$ has no
smaller nonempty open on which it could split. This is the practical
reason to allow more than Zariski covers.

\subsection{What stability and composition say}\label{subsec:coverrules}
\emph{Stable under base change} means the following concrete operation
preserves the covering property. If $\{S_i\to S\}$ is a cover and
$f:T\to S$ is \emph{any} $B$-morphism, form
\[
\begin{tikzcd}[column sep=large,row sep=large]
T\times_S S_i \arrow[r] \arrow[d] &S_i \arrow[d]\\
T \arrow[r,"f"'] &S.
\end{tikzcd}
\]
The conclusion is that $\{T\times_S S_i\to T\}_i$ is again an fppf
cover. Flatness and local finite presentation survive this pullback,
and the pulled-back family is jointly surjective. Geometrically, a
point of $T$ has image in some $S_i\to S$ and has a point above it in
the corresponding fiber product, allowing a residue-field extension
if necessary. Thus the local pieces still cover after changing base.

\emph{Composition of covers} means exactly covering the covering
pieces. If $\{S_i\to S\}_i$ is a cover and for each $i$ we have a
cover $\{U_{ij}\to S_i\}_{j\in J_i}$, then the composite family
\[
 \{U_{ij}\longrightarrow S_i\longrightarrow S\}_{i,j}
\]
covers $S$. Composites remain flat and locally of finite presentation;
joint surjectivity follows in two steps. The index set $J_i$ is allowed
to depend on $i$.

The exact reference is Lemma~34.7.3, \BGtag{021O}: part~(1) concerns
isomorphisms, part~(2) composition, and part~(3) base change. In
contrast, Lemma~34.7.2, \BGtag{021N}, compares types of covers: it says
that syntomic covers, and hence smooth, \emph{\'etale}, and Zariski
covers, are fppf. It does \emph{not} establish the preceding two
stability statements. In particular the inclusion is not a statement
that every fppf cover is smooth or \emph{\'etale}.

These verifications are needed even though our initial covers are
fppf: pulling back a torsor will produce a \emph{new} covering family,
and gluing locally trivial objects will produce a cover of a cover.
The lemma ensures those new families are still allowed.

\subsection{Locality and the use of fpqc theorems}
``After this cover'' means after pullback along every $S_i\to S$.
For a scheme $P\to S$, write $P_{S_i}=P\times_S S_i$. Thus
``$P$ is locally isomorphic to $Q$'' means $P_{S_i}\simeq Q_{S_i}$
for some cover. It need not mean an isomorphism on Zariski open
subsets of $S$.

We shall occasionally quote theorems stated for \emph{fpqc} covers,
a broader class of flat covers that permits maps not locally of finite
presentation. In the covering-family definition, every affine open
in the target must equal the union of the images of finitely many
affine opens in members of the family. Affine opens are quasi-compact:
a quasi-compact space
is one in which every open cover has a finite subcover. This condition
provides the finiteness needed in descent without imposing finitely
many algebra generators. This is Definition~34.9.1, \BGtag{022B}.

The fact we use is Lemma~34.9.7, \BGtag{022C}: every fppf cover is
an fpqc cover. Consequently an fpqc descent theorem applies to each
of \emph{our fppf covers}. We are not changing topology in the proof.
Just as a theorem for all groups can be applied to finite groups, a
theorem for the larger class can be used on this subclass.

Two further phrases must be distinguished. A property is
\emph{preserved by base change} if a map having it still has it after
any change of base. It is \emph{detected by a cover} if knowing it
for all pulled-back maps implies it for the given original map.
For example, if a given map $P\to Q$ becomes an isomorphism after an
fppf cover of $S$, a locality theorem says it was already an
isomorphism. The original $P,Q$ must already exist for this assertion.
``Fpqc-local on the target'' packages preservation along such covers
and detection on them; it does not by itself assert preservation
under \emph{arbitrary}, possibly nonflat, base change.

\subsection{Presheaves and sheaves: what exactly gets glued?}
A \emph{presheaf of sets} assigns a set $F(T)$ to each object $T$ and
a restriction function $F(T)\to F(T')$ to each arrow $u:T'\to T$.
Write $u^*s$ for the restriction of $s$. The laws are
\[
 \id^*s=s,\qquad v^*(u^*s)=(uv)^*s.
\]
This is the same reversal of arrows as for the functor of points.
It does not yet assert that locally defined elements glue.

A presheaf is a \emph{sheaf} for our covering rule if, whenever
$\{T_i\to T\}$ is a cover and $s_i\in F(T_i)$ have equal
restrictions to every $T_i\times_T T_j$, there exists exactly one
$s\in F(T)$ restricting to all $s_i$. For open covers this is the
familiar agreement-on-intersections rule. Fiber products replace
intersections because the new covering pieces need not be subsets.

For example, $F(T)$ might be functions on $T$, or maps from $T$ to
a fixed scheme, or isomorphisms between two objects pulled back to $T$.
In the last case the \emph{elements being glued are maps}. Later the
second stack condition will glue the \emph{objects themselves}; that
is a separate problem.

A map of presheaves is a family of functions on $F(T)$ commuting
with restrictions, just as in Section~\ref{sec:points}. A presheaf
$F$ is \emph{represented by a scheme $X$} if there are compatible
bijections $F(T)\simeq\Hom(T,X)$ for every $T$ in the base category.
This says that the rule $F$ comes from an actual geometric object $X$.
Not every rule comes with such an object.

\BGcheckpoint{\begin{itemize}
\item Defined: site, fppf cover, presheaf, sheaf, and representability.
\item Verified: fppf covers survive pullback and composition.
\item Check: distinguish an fppf cover from an arbitrary arrow of the site.
\item Check: distinguish gluing maps from gluing objects.
\end{itemize}}

\section{Torsors: groups without a chosen origin}\label{sec:torsors}

\subsection{Actions, sections, and triviality}
For an ordinary group $H$, a nonempty set on which $H$ acts freely
and transitively looks like $H$ after one chooses an element to be its
origin. Without that choice it is a torsor, not a group with a specified
identity. Our scheme definition makes this description work in families.

Given a scheme $P\to S$, an action of $G_S$ is an $S$-morphism
$a_P:G_S\times_S P\to P$ satisfying
$a_P(e,p)=p$ and $a_P(g,a_P(h,p))=a_P(gh,p)$ on all test points.
Equivalently these are equations of scheme morphisms. An equivariant
map respects this action: applying $g$ and then the map is the same
as applying the map and then $g$. See Definition~39.10.1,
\BGtag{022Z}.

A \emph{section} of $\pi_P:P\to S$ is a morphism $s:S\to P$ with
$\pi_Ps=\id_S$. It chooses an element over every base point in a way
that is a morphism of schemes, not an arbitrary set-theoretic choice.
For a double cover with sheet swap, such a section labels its sheet
as $0$ and the other sheet as $1$. A torsor need not have a global
section, but it must have such descriptions after allowed base changes.

\subsection{The definition with its local isomorphisms}
\begin{Def}\label{def:torsor}
A left $G_S$-torsor is a scheme $\pi_P:P\to S$ with an action
\[
 a_P:G_S\times_S P\longrightarrow P
\]
and an fppf cover $\{U_a\to S\}$ on which there are equivariant
isomorphisms $P_{U_a}\simeq G_{U_a}$. On $G_{U_a}$ the action is left
multiplication. The action must satisfy the unit and associativity laws.
\end{Def}
Our convention agrees with Definition~39.11.3, \BGtag{049A}.
Trivializations are required to exist; they are not extra chosen data in
the definition of the object.

The associated \emph{torsor identity} is the map
\begin{equation}\label{eq:theta}
\begin{aligned}
 \Theta_P:G_S\times_S P&\longrightarrow P\times_S P,\\
 (g,p)&\longmapsto(g\cdot p,p).
\end{aligned}
\end{equation}
It is an isomorphism. On a trivialization its inverse is
$(p_1,p_2)\mapsto(p_1p_2^{-1},p_2)$, and being an isomorphism is
fpqc-local on the base (Lemma~35.23.19, \BGtag{02L4}).

In particular, this asserts simple transitivity on sections after any
test-scheme base change whenever a section is available. It is stronger
than a statement only about geometric points.

\begin{Ex}[The torsor identity on four copies]
Take $G=(\bZ/2\bZ)_B$ and the trivial torsor $P=S_0\amalg S_1$,
where the subscripts label the two copies of $S$. Both
$G_S\times_S P$ and $P\times_S P$ consist of four copies of $S$.
The map $\Theta_P$ acts on their labels as follows:
\[
\begin{array}{c|c}
 \text{source }(g,p)&\text{target }(g+p,p)\\ \hline
 (0,0)&(0,0)\\
 (1,0)&(1,0)\\
 (0,1)&(1,1)\\
 (1,1)&(0,1)
\end{array}
\]
On each labeled copy it is the identity map of $S$. Thus this table
describes an isomorphism of schemes, not merely a bijection between
four abstract points. Given target labels $(a,b)$, the unique group
label carrying $b$ to $a$ is $g=a-b$ in $\bZ/2\bZ$. The inverse
sends $(a,b)$ to $(a-b,b)$ and again preserves the base. This is what
``there is a unique group element taking the second point to the
first'' looks like as a scheme morphism.
\end{Ex}

\begin{Prop}\label{prop:torsorcriterion}
Under the standing assumptions on $G$, Definition~\ref{def:torsor} is
equivalent to requiring that $P\to S$ be surjective, flat, and locally
of finite presentation and that $\Theta_P$ be an isomorphism.
\end{Prop}
\begin{proof}
Local triviality proves the identity as above and transfers flatness and
local finite presentation from $G_S\to S$ to $P\to S$. These properties
are fpqc-local on the base (Lemmas~35.23.17 and~35.23.13,
\BGtag{02L2} and \BGtag{02KY}). Surjectivity follows on the same
trivializing cover, since $G_S$ has a unit section.

Conversely, use the singleton fppf cover $P\to S$. The pulled-back
object $P\times_S P\to P$, with projection to its second factor, is
equivariantly isomorphic via $\Theta_P$ to $G_S\times_S P\to P$.
It is therefore trivial on this cover.
\end{proof}

More generally, a section $s:S\to P$ trivializes a torsor by
$g\mapsto g\cdot s$. Its inverse sends $p$ to the unique $g$ with
$p=g\cdot s(\pi_P(p))$, furnished by $\Theta_P^{-1}$. This is the
section criterion of Lemma~39.11.2, \BGtag{0499}.

\subsection{Sheaf torsors and representability}
The proof below can be followed using schemes as torsors throughout.
This subsection explains a different definition often encountered in
references. It separates a rule for local sections from the question
whether a scheme realizes that rule.
The representability paragraph previews the gluing theorem of
Section~\ref{sec:effective}; its ``cocycle'' means the consistency
equation for three local identifications, defined fully in
Section~\ref{sec:stackdef}.

For a scheme $X\to S$, write $h_X(T)=\Hom_S(T,X)$. A topology is
\emph{subcanonical} if each such representable functor is a sheaf.
The fppf topology is subcanonical: schemes satisfy the stronger fpqc
sheaf condition (Lemma~35.13.7, \BGtag{023Q}).

A torsor under the sheaf of groups $h_{G_S}$ is a sheaf of sets
$\cP$ on $(\Sch/S)_{\mathrm{fppf}}$, with an action, locally isomorphic
to $h_{G_S}$ with its regular action. No scheme is part of this definition.
For an already represented sheaf $h_P$, the two torsor definitions agree
(Lemma~39.11.4, \BGtag{049B}).

For affine $G$, \emph{every} such sheaf torsor is represented by an affine
$S$-scheme. Here is the representability argument and its precise input.
Choose local trivializations by affine $U_a$-schemes $G_{U_a}$. Yoneda
turns their sheaf transition maps into scheme isomorphisms satisfying a
cocycle. Effective descent for affine morphisms, stated in
Theorem~\ref{th:affinedescent}, produces an affine $S$-scheme $P$.
The local identifications $h_P|_{U_a}\simeq\cP|_{U_a}$ glue to
$h_P\simeq\cP$ by the sheaf axiom. Thus representability uses affine
descent, not the claim that $BG$ is a stack; compare Remark~39.11.5,
\BGtag{049C}. One can also allow geometric objects called algebraic
spaces, which are needed for the comparison with curves later.
No knowledge of them is needed in the torsor proof: affineness lets
us stay inside schemes.

\begin{Ex}[The first double cover]\label{ex:trivialdouble}
Let $G=(\bZ/2\bZ)_B=B\amalg B$. Then
\[
 G_S\times_S S\simeq G_S=S\amalg S\to S
\]
is the trivial torsor. The nonidentity element exchanges the two sheets.
This group scheme is finite \emph{\'etale} over every base. Its sections
over $T$ are locally constant maps from $T$ to $\bZ/2\bZ$.

On an affine base $\Spec A$,
\[
 \mu_2=\Spec A[z]/(z^2-1).
\]
If $2$ is invertible, the Chinese remainder theorem identifies it with
$\Spec(A\times A)$. Over a nonempty $\bF_2$-scheme it has a nonzero
square-zero coordinate $z-1$, whereas the constant group has two disjoint
copies of the base and is \emph{\'etale}. The two groups are therefore
different in characteristic $2$. This distinction is relative to the
base, which itself might be nonreduced. Over a characteristic-two
field-valued base point, their fibers are respectively
$\Spec k[\epsilon]/(\epsilon^2)$ and
$\Spec k\amalg\Spec k$, so they cannot be isomorphic over that base.
\end{Ex}

\begin{Ej}[The torsor identity alone]\label{ex:empty}
Let $S\ne\varnothing$. Show that $P=\varnothing$ satisfies
$G_S\times_S P\simeq P\times_S P$ but is not a torsor. Identify the
missing condition.
\end{Ej}
\begin{Ej}[Sections and triviality]\label{ex:section}
Prove that a torsor $P\to S$ is trivial if and only if it has a section.
Describe how changing a section changes the trivialization.
\end{Ej}
\begin{Ej}[Two groups of order two]\label{ex:mu}
Compare the coordinate algebras of $(\bZ/2\bZ)_{\Spec k}$ and
$\mu_{2,k}$ when $\Char k=2$. What are their $k$-points? Explain why a
comparison only of orders or ranks cannot identify the group schemes.
\end{Ej}

\subsection{The topology choice revisited}
The single fppf cover $\Spec\bC\to\Spec\bR$ already showed why
Zariski local triviality is too restrictive. For the constant group of
order two, the torsors are finite \emph{\'etale} double covers; this
means in particular that they split on \emph{\'etale} covers, a class
of algebraic local coordinate changes broader than open immersions.
The precise splitting theorem is used in Section~\ref{sec:double}.
All these covers are fppf, so we can keep one covering rule throughout.

For reference, if $G\to B$ is smooth, an fppf torsor is smooth and
surjective over its base: smoothness is detected on the trivializing
cover (Lemma~35.23.29, \BGtag{02VL}). Smooth surjective morphisms have
\emph{\'etale} local sections (Lemma~37.38.6, \BGtag{055U}), so
such torsors are already \emph{\'etale}-locally trivial. The fppf
rule also accommodates nonsmooth groups such as $\mu_2$ in
characteristic $2$. Under our flatness and presentation hypotheses,
fpqc-local triviality gives the same torsors (Lemma~95.15.4,
\BGtag{0CQJ}). These comparisons are background; the proof uses fppf
covers and the explicitly stated descent theorems only.

\BGcheckpoint{\begin{itemize}
\item Defined: action, equivariance, section, torsor, and local triviality.
\item Verified: a section trivializes a torsor; the torsor identity
expresses simple transitivity.
\item Check: explain why a double cover can be locally trivial without
a global choice of sheet.
\end{itemize}}

\begin{Th}[Target theorem]\label{th:target}
Let $B$ be a scheme and let $G\to B$ satisfy
Definition~\ref{def:standing}. The category of scheme-represented
fppf $G$-torsors, with equivariant arrows over varying $B$-schemes,
is a stack in groupoids on $(\Sch/B)_{\mathrm{fppf}}$.
It realizes
\[
 BG=[B/G]=[\operatorname{pt}/G],
\]
where $\operatorname{pt}$ is the relative point $B$.
\end{Th}
This is the assertion to be proved, not an assumption. The next
sections first construct the category and then verify each of the
properties in the roadmap.

\section{The meaning of the quotient notation}\label{sec:quotient}

Before constructing the total category, distinguish four meanings of a
quotient by a group.

First consider an ordinary group $H$. Make a category with one object,
called $*$, and one arrow $*\to *$ for each element of $H$. Compose
arrows using multiplication in $H$; the identity element gives the
identity arrow, and $h^{-1}$ gives the inverse of the arrow $h$.
A category whose arrows are all invertible is called a \emph{groupoid}.
An invertible arrow from an object to itself is an \emph{automorphism}.
Thus this one-object groupoid remembers the whole group of symmetries,
although its set of objects has only one member.

For a $B$-scheme $S$, we can do this with the ordinary group
$G(S)=\Hom_B(S,G)$. The important point is that changing $S$ changes
the group of sections. We must retain this dependence on $S$, not just
choose one group once and for all.

\emph{The action groupoid.} The trivial action on the relative point is
encoded by the diagram $G\rightrightarrows B$: both arrows are the
structure map, while multiplication in $G$ specifies composition.
This is called an \emph{internal groupoid} because its objects and
arrows are themselves specified by schemes. For each $B$-scheme $S$,
it gives the one-object groupoid with automorphism group $G(S)$ just
described.
It retains the stabilizer, meaning the group of symmetries of the
object; the set of orbits would be a singleton.

\emph{The naive quotient prestack.} Varying $S$ gives a groupoid-valued
presheaf $\cB_0$ with one object over $S$ and automorphisms $G(S)$.
Pullback of arrows is restriction of sections. Because $G$ is a sheaf,
arrows already satisfy sheaf descent; this is the meaning of
\emph{prestack} used here. Object descent is not yet required.

\emph{Stackification.} A stackification of a prestack is a stack
receiving a morphism from it, universal for morphisms to stacks, up to
the usual natural isomorphisms. For the action groupoid this construction
is the definition of the quotient stack; see Definition~78.20.1,
\BGtag{044Q}. Informally it adds objects represented by compatible local
objects, and identifies descriptions related by refinement.

\emph{Torsors.} The unique object of $(\cB_0)_S$ corresponds to the
trivial torsor $G_S$. This means that the one object is represented
geometrically by the trivial family, and its arrows become symmetries
of that family.

\emph{Convention check.} The next calculation only keeps track of
multiplication order; for the running example $\bZ/2\bZ$, the group is
commutative and every element equals its inverse, so this distinction
disappears. A left-equivariant automorphism of $G_S$ has the
form $R_h(x)=xh$. Since
\[
 R_k\circ R_h=R_{hk},
\]
$\Aut_{G_S}(G_S)$ is naturally $G(S)^{\mathrm{op}}$ via $h\mapsto R_h$.
The functor from the one-object groupoid with composition
$g_2\circ g_1=g_2g_1$ sends $g$ to $R_{g^{-1}}$; indeed
\[
 R_{g_2^{-1}}\circ R_{g_1^{-1}}
 =R_{(g_2g_1)^{-1}}.
\]
This makes the left/right convention explicit.

Every torsor is locally a trivial one, and its local identifications are
the transition automorphisms of those trivial objects. Once the stack
axioms are proved below, this identifies the resulting stack with the
stackification of $\cB_0$. The canonical equivalence is
Lemma~95.15.4, \BGtag{0CQJ}; it is a comparison with the notation, not an
input into our proof of descent.

For $G=\bZ/2\bZ$, the naive fiber contains just a labeled pair of
sheets, while stackification admits all locally split double covers.
Section~\ref{sec:double} gives explicit nonsplit examples.

\begin{Ej}[Opposite-group check]\label{ex:opposite}
For a possibly noncommutative $G$, verify the two displayed composition
formulas. Explain why $g\mapsto R_g$ is not generally a functor from
the one-object groupoid with its usual composition, and why
$g\mapsto R_{g^{-1}}$ is.
\end{Ej}

\BGcheckpoint{\begin{itemize}
\item Distinguished: the one-object groupoid, the naive prestack,
its stackification, and the torsor description.
\item Check: explain why a nonsplit double cover is missing from the
naive fiber.
\end{itemize}}

\section{Checking that BG is an ordinary category}\label{sec:total}

\subsection{Objects, arrows, and equivariance}
We return to Definition~\ref{def:BGdata}. Its objects are the pairs
$(S,P)$, with $P\to S$ a $G_S$-torsor. We now understand all the
ingredients. We spell out equivariance for its arrows, then supply
identities and composition and verify the category axioms.

For $f:S\to T$, write
\[
 \alpha_f=\id_G\times f:G_S\to G_T.
\]
A prospective arrow $(S,P)\to(T,Q)$ is a pair $(f,\widetilde f)$ with
$f$ a $B$-morphism and $\widetilde f:P\to Q$ fitting into
\[
\begin{tikzcd}[column sep=large,row sep=large]
P \arrow[r,"\widetilde f"] \arrow[d,"\pi_P"'] & Q \arrow[d,"\pi_Q"]\\
S \arrow[r,"f"'] & T
\end{tikzcd}
\]
The required equivariance means that the following diagram commutes:
\[
\begin{tikzcd}[column sep=small,row sep=large]
G_S\times_S P \arrow[r,"a_P"]
  \arrow[d,"\alpha_f\times\widetilde f"'] & P \arrow[d,"\widetilde f"]\\
G_T\times_T Q \arrow[r,"a_Q"'] & Q
\end{tikzcd}
\]
In test-point notation this is
$\widetilde f(g\cdot p)=\alpha_f(g)\cdot\widetilde f(p)$.
The factor $\alpha_f$ is necessary because the groups act over different
bases.

\subsection{Category axioms}
Define
\[
 \id_{(S,P)}=(\id_S,\id_P).
\]
Both required diagrams commute for this pair. For successive arrows
$(f,\widetilde f)$ and $(g,\widetilde g)$, define
\[
 (g,\widetilde g)(f,\widetilde f)
 =(gf,\widetilde g\widetilde f).
\]
The composite lies over $gf$ since
\[
 \pi_R\widetilde g\widetilde f
 =g\pi_Q\widetilde f=gf\pi_P.
\]
Its equivariance follows from $\alpha_{gf}=\alpha_g\alpha_f$:
\[
\begin{aligned}
 \widetilde g\widetilde f(h\cdot p)
 &=\widetilde g(\alpha_f(h)\cdot\widetilde f(p))\\
 &=\alpha_{gf}(h)\cdot\widetilde g\widetilde f(p).
\end{aligned}
\]
For three arrows, the base components of the two parenthesizations agree
by associativity in $\Sch/B$, and the total-space components agree by
associativity of scheme morphisms. Thus composition is associative.
Composing on the left with $(\id_T,\id_Q)$, or on the right with
$(\id_S,\id_P)$, leaves both components of $(f,\widetilde f)$ unchanged.
This proves both identity laws.

\BGcheckpoint{\begin{itemize}
\item Defined: objects, arrows, identities, and composition.
\item Verified: equivariance of composites, associativity, and both
identity laws.
\item Result: $BG$ is an ordinary category.
\end{itemize}}

For double covers, equivariance says
$\widetilde f\sigma_P=\sigma_Q\widetilde f$. Identities and composites
commute with the involutions by the equations just checked. We do not
define arrows to be arbitrary maps of degree-two covers with the actions
forgotten.

\begin{Ex}[Moving the base and exchanging the sheets]
Let $B=\Spec k$, $S=\bA^1_k$, and $P=S_0\amalg S_1$ with its
sheet-swap action. The following square defines a total-category
arrow $(S,P)\to(S,P)$:
\[
\begin{tikzcd}[column sep=large,row sep=large]
 S_0\amalg S_1 \arrow[r,"\widetilde f"] \arrow[d,"\pi"']
 &S_0\amalg S_1 \arrow[d,"\pi"]\\
 S \arrow[r,"f"']&S.
\end{tikzcd}
\]
Here $f$ is the squaring map, and on labeled test points
\[
 f(t)=t^2,\qquad \widetilde f(t,a)=(t^2,a+1).
\]
Each component map is the scheme morphism defined by
$k[t]\to k[t]$, $t\mapsto t^2$. The equation expressing
equivariance is
\[
 \widetilde f(t,a+1)=(t^2,a)=\sigma\widetilde f(t,a).
\]
This arrow is not invertible: an inverse pair would require an inverse
to $f$, but the image of its ring map does not contain $t$. In contrast,
$\sigma(t,a)=(t,a+1)$ together with $\id_S$ defines an invertible
arrow whose square is the identity. Composing the displayed arrow
with itself gives $(t,a)\mapsto(t^4,a)$: the two sheet exchanges
cancel, but the composite base map is not the identity. The two components of a
total-category arrow carry different information.
\end{Ex}

\begin{Ej}[Category axioms]\label{ex:axioms}
Write the complete verification of the identity laws and associativity
for three pairs $(f,\widetilde f)$. Explain why defining arrows as actual
maps $P\to Q$ avoids having to insert associativity isomorphisms of
fiber products in the definition of composition.
\end{Ej}
\begin{Ej}[All arrows between trivial torsors]\label{ex:arrows}
Fix $f:S\to T$. Show that an equivariant map $G_S\to G_T$ over $f$
has the form
\[
 (x,s)\longmapsto(xh(s),f(s)),
 \qquad h\in G(S).
\]
If a second arrow over $g:T\to U$ is determined by $k\in G(T)$,
compute the section determining the composite.
\end{Ej}

\section{The projection to schemes}\label{sec:projection}

The definition of a category over a base now applies to our constructed
category. The data already include base objects and base arrows; the
remaining obligation is to check functoriality. Define
\[
 p:BG\longrightarrow\Sch/B,
 \qquad p(S,P)=S,
\]
and $p(f,\widetilde f)=f$. We have
\[
 p(\id_{(S,P)})=\id_S
\]
and
\[
\begin{aligned}
 p((g,\widetilde g)(f,\widetilde f))
 &=gf\\
 &=p(g,\widetilde g)p(f,\widetilde f).
\end{aligned}
\]
Thus this is a functor. It forgets the torsor but retains the base map of
each arrow. The terminology ``over $S$'' and ``over $f$'' now refers
precisely to this functor.

\BGcheckpoint{\begin{itemize}
\item Defined: $p(S,P)=S$ and $p(f,\widetilde f)=f$.
\item Verified: preservation of identities and composition.
\item Result: $p:BG\to\Sch/B$ is a functor.
\end{itemize}}

For a double cover, $p$ remembers $S$ and forgets both sheets and the
involution. In particular, forgetting the involution is not the same
operation as taking isomorphism classes of covers.

\begin{Ej}[An arrow over a nonidentity endomorphism]\label{ex:projection}
Assume $B\ne\varnothing$. Take $S=\bA^1_B$ with its endomorphism $f$
given by $t\mapsto t^2$.
Use trivial torsors to construct an arrow $(S,G_S)\to(S,G_S)$ over $f$.
Why is it not a morphism in the fiber $(BG)_S$?
\end{Ej}

\section{Fiber categories and pullback torsors}\label{sec:pullback}

\subsection{The fiber, without yet asserting invertibility}
By definition, $(BG)_S$ has $G_S$-torsors as objects and equivariant
$S$-maps as arrows. In other words, the arrows have base component
$\id_S$. At this point we do not build invertibility into their
definition. It will be proved in Section~\ref{sec:groupoids}.

\subsection{Keep the types of objects separate}
The following similar-looking symbols describe different kinds of
mathematical objects:
\begin{itemize}
\item $G_S$ is a \emph{scheme over $S$}, with a group law.
\item $(BG)_S$ is a \emph{category}: its objects are torsors over
$S$, and its arrows are equivariant $S$-maps between them.
\item $G_S(T)=\Hom_S(T,G_S)$ is a \emph{set with a group law},
once a particular morphism $T\to S$ has been specified.
\item $P_T=P\times_S T$ is a \emph{scheme over $T$}, obtained
from a particular torsor $P\to S$; it is one object of $(BG)_T$.
\end{itemize}
For the trivial double cover, $G_S$ and $P$ can be the same
underlying $S$-scheme. Neither is the category $(BG)_S$: that
category also includes other torsors over $S$ and the maps between
them. The subscript on $(BG)_S$ means ``keep the base fixed at $S$,''
whereas the parentheses in $G_S(T)$ mean ``take maps from $T$.''

\subsection{Base change of the geometric data}
Given $f:T\to S$ and a torsor $P\to S$, we must produce an object over
$T$. Set
\[
 f^*P=P\times_S T,\qquad G_T=G_S\times_S T.
\]
The base change of the action is
\[
 g\cdot(p,t)=(\alpha_f(g)\cdot p,t).
\]
The unit and associativity equations hold because their diagrams are
base changes of the action diagrams on $P$. If $\{U_a\to S\}$
trivializes $P$, then $\{U_a\times_S T\to T\}$ is an fppf cover, and
on it $f^*P$ is the regular $G_T$-torsor. Consequently $f^*P\to T$ is
a torsor. The morphism $\Theta_{f^*P}$ is also the base change of
$\Theta_P$, so it is an isomorphism.

Projection gives an equivariant arrow
\begin{equation}\label{eq:pullbackarrow}
 (f,\pr_P):(T,f^*P)\longrightarrow(S,P).
\end{equation}
We have constructed a lift of $f$; it remains to prove its universal
property in the total category.

\BGcheckpoint{\begin{itemize}
\item Constructed: $f^*P$, its action, and its projection to $P$.
\item Verified: local triviality and equivariance.
\item Still to check: the cartesian universal property.
\end{itemize}}

For $\bZ/2\bZ$, pullback gives $P\times_S T\to T$ with involution
$(p,t)\mapsto(\sigma_P(p),t)$. Pulling back $S\amalg S\to S$ yields
$T\amalg T\to T$.

\begin{Ej}[A pullback computation]\label{ex:pullback}
Let $k$ have characteristic different from $2$, and consider
$P\to S=\Spec k[t,t^{-1}]$ given by $u^2=t$.
Pull it back by $T=\Spec k[v,v^{-1}]\to S$, $t=v^2$.
Write its coordinate algebra and exhibit its two sheets and involution.
\end{Ej}

\section{Cartesian morphisms and the fibration}\label{sec:cartesian}

The pullback torsor has already been constructed. Its extra property
is this: a map into the original torsor can be viewed as a map into the
pulled-back torsor whenever its base map has been specified through the
new base. There is exactly one way to do so. The word \emph{cartesian}
packages this existence-and-uniqueness assertion; it does not mean that
every object mapping into the pullback is isomorphic to it.

\begin{Def}\label{def:cartesian}
An arrow $u:x\to y$ in a category over $\cC$, lying over $f:T\to S$,
is \emph{cartesian} if the following holds. Given $z$ over $V$, an arrow
$v:z\to y$ over $h:V\to S$, and $k:V\to T$ with $h=fk$, there is a
unique arrow $w:z\to x$ over $k$ such that $uw=v$.
A \emph{fibered category} is a category over $\cC$ admitting a
cartesian lift of every $f:T\to S$ with every prescribed target over $S$.
\end{Def}
The Stacks Project uses the term \emph{strongly cartesian} for this
property, which allows any factorization $h=fk$
(Definition~4.33.1, \BGtag{02XK}). Its fibration definition requires
these lifts (Definition~4.33.5, \BGtag{02XM}).

The data already constructed in \eqref{eq:pullbackarrow} have the right
source, target, and base map. We now verify the quantified factorization
condition. Let $(V,R)$ be any object, let
$(h,\widetilde h):(V,R)\to(S,P)$ be any arrow, and suppose $h=fk$.
The competing diagrams are
\[
\begin{tikzcd}[column sep=small,row sep=large]
R \arrow[rr,"\widetilde h"] \arrow[dr,dashed,"\lambda"'] && P\\
& P\times_S T \arrow[ur,"\pr_P"'] &
\end{tikzcd}
\]
over
\[
\begin{tikzcd}[column sep=small,row sep=large]
V \arrow[rr,"h"] \arrow[dr,"k"'] && S\\
& T \arrow[ur,"f"'] &
\end{tikzcd}
\]

The two maps $\widetilde h:R\to P$ and $k\pi_R:R\to T$ satisfy
\[
 \pi_P\widetilde h=h\pi_R=fk\pi_R.
\]
The universal property of $P\times_S T$, stated in
Section~\ref{sec:categories}, therefore gives a unique map
\[
\begin{aligned}
 \lambda:R&\longrightarrow P\times_S T,\\
 r&\longmapsto(\widetilde h(r),k\pi_R(r)).
\end{aligned}
\]
It lies over $k$, and $\pr_P\lambda=\widetilde h$. For $g\in G_V$,
\[
\begin{aligned}
 \lambda(g\cdot r)
 &=(\alpha_h(g)\cdot\widetilde h(r),k\pi_R(r))\\
 &=\alpha_k(g)\cdot\lambda(r),
\end{aligned}
\]
because $\alpha_h=\alpha_f\alpha_k$. Thus $(k,\lambda)$ is an arrow of
$BG$. Any other factorization must have the same projections to $P$
and $T$, so the fiber-product uniqueness clause makes it equal to
$(k,\lambda)$. This proves cartesianness.

In the double-cover picture, the map $\lambda$ records two things:
which point of the old cover $P$ a point of $R$ maps to, and which
point of the new base $T$ it is meant to lie above. The equation
$h=fk$ says these two records describe the same point of $S$.
It makes the pair a point of the fiber product. The construction gives
a map $R\to P\times_S T$, not an assertion that $R$ equals that
fiber product.

\subsection{Unique pullbacks and their coherence}
Suppose $u:x\to y$ and $u':x'\to y$ are cartesian lifts of the same
$f:T\to S$. Applying their universal properties gives vertical arrows
$a:x\to x'$ and $b:x'\to x$ with $u'a=u$ and $ub=u'$.
Both $ba$ and $\id_x$ are factorizations of $u$ through $u$ over
$\id_T$, so $ba=\id_x$; likewise $ab=\id_{x'}$. The isomorphism $a$
is unique \emph{subject to compatibility with $u,u'$}. This does not
forbid other automorphisms of the underlying torsor.

For $T\xrightarrow{f}S\xrightarrow{g}U$ and a torsor $P\to U$,
the canonical comparison is
\[
\begin{aligned}
 P\times_U T&\xrightarrow{\sim}(P\times_U S)\times_S T,\\
 (p,t)&\longmapsto((p,f(t)),t).
\end{aligned}
\]
Thus $(gf)^*P\simeq f^*(g^*P)$. Projection also gives
$(\id_S)^*P\simeq P$. The comparison maps are equivariant and respect
the maps to $P$. For three successive base changes, all ways of inserting
these comparisons agree: they have the same projections to each factor,
so uniqueness in the fiber-product property proves coherence.
The general categorical version is Lemma~4.33.7, \BGtag{02XO}.

\BGcheckpoint{\begin{itemize}
\item Verified: a competing arrow factors uniquely through the pullback
arrow over the prescribed base map.
\item Result: $BG\to\Sch/B$ is fibered.
\item Check: write both projections of the factorization.
\end{itemize}}

For double covers, the factorization sends $r$ to its image in the
original cover together with its prescribed point of the new base.
Successive base changes retain the same two sheets and give canonically
isomorphic double covers.

\begin{Ej}[Complete the cartesian diagram]\label{ex:cartesian}
In the competing diagram above, prescribe $R$, $\widetilde h$, and $k$.
Write both projections of $\lambda$. Prove directly that equivariance
uses $\alpha_h=\alpha_f\alpha_k$, and that no other arrow over $k$ can
factor $(h,\widetilde h)$ through \eqref{eq:pullbackarrow}.
\end{Ej}
\begin{Ej}[Uniqueness with a qualification]\label{ex:unique}
A trivial double cover has a nontrivial deck automorphism. Explain why
this does not contradict uniqueness up to unique isomorphism of a
cartesian lift. Which compatibility equation does the deck automorphism
usually fail?
\end{Ej}

\section{Why the fibers are groupoids}\label{sec:groupoids}

\begin{Def}
A groupoid is a category in which every arrow has an inverse. A
\emph{category fibered in groupoids} is a fibered category whose fibers
are groupoids.
\end{Def}

We already have a fibered category, and $(BG)_S$ was defined using all
equivariant $S$-maps. The remaining task is to prove their invertibility.

\begin{Lem}[Equivariant maps of torsors]\label{lem:torsormap}
If $P$ and $Q$ are $G_S$-torsors, every equivariant $S$-map $u:P\to Q$
is an isomorphism, whose inverse is equivariant.
\end{Lem}
\begin{proof}
Take trivializing covers for both torsors and form their fiber products;
this gives a common fppf trivializing cover $\{U_a\to S\}$.
On each $U_a$, the map becomes an equivariant map $G_{U_a}\to G_{U_a}$.
If $h=u(e)$, equivariance gives $u(x)=xh$. The map $x\mapsto xh^{-1}$
is its inverse. Consequently $u$ is an isomorphism after the cover and
therefore globally by \BGtag{02L4}. If $q=u(p)$, the equation
$u(gp)=gq$ implies $u^{-1}(gq)=gu^{-1}(q)$, proving equivariance of the
inverse.
\end{proof}

\BGcheckpoint{\begin{itemize}
\item Proved: every equivariant map of torsors over the same base is
invertible.
\item Result: every fiber $(BG)_S$ is a groupoid.
\item Together with cartesian lifts: $BG\to\Sch/B$ is a category
fibered in groupoids.
\end{itemize}}

\subsection{The equivalent description of total arrows}
We can now justify an often-used alternative definition. Given
$(f,\widetilde f):(S,P)\to(T,Q)$, the commutative square induces
\[
\begin{aligned}
 \widehat f:P&\longrightarrow f^*Q=Q\times_T S,\\
 p&\longmapsto(\widetilde f(p),\pi_P(p)).
\end{aligned}
\]
It is an equivariant $S$-map between torsors, hence an isomorphism by
Lemma~\ref{lem:torsormap}. Conversely, such an isomorphism followed by
projection to $Q$ yields $\widetilde f$. Therefore total arrows are
equivalently base maps $f$ equipped with $P\simeq f^*Q$; see also
Examples of Stacks, Section~95.14.8, \BGtag{04UR}.

Every total arrow is consequently cartesian: it is the composite of a
vertical isomorphism with a cartesian pullback arrow. This also verifies
the alternative convention for a category fibered in groupoids in which
all arrows are required to be cartesian. Notice that $P\to Q$ itself
need not be an isomorphism when the base map is not an isomorphism.

For the constant group of order two, an arrow in the fiber is locally
$x\mapsto x+\varepsilon$ with $\varepsilon\in\bZ/2\bZ$. It is either
the identity or the sheet swap on each connected component. This is the
concrete form of the lemma.

\begin{Ej}[Invertibility without defining it]\label{ex:invertible}
Prove Lemma~\ref{lem:torsormap}, including the construction of a common
trivializing cover. Why is checking bijectivity only on geometric fibers
not a substitute for the scheme-theoretic proof? Give a total arrow
whose map on total spaces is not an isomorphism.
\end{Ej}

\section{The definition of a stack}\label{sec:stackdef}

The available structure is now a category fibered in groupoids. A choice
of cartesian lifts is called a \emph{cleavage}; it makes notation such as
$f^*x$ convenient. Its canonical comparison isomorphisms will always be
inserted when comparing iterated pullbacks.

\begin{Def}[The two stack conditions]\label{def:stack}
A category fibered in groupoids $\cF$ on our site is a stack if:
\begin{enumerate}[label=(\arabic*)]
\item Isomorphisms form a sheaf. For $x,y\in\cF_S$, define
\[
 T\longmapsto\Isom_{\cF_T}(x_T,y_T).
\]
This presheaf must be a sheaf for fppf covers of $S$-schemes.
\item For every fppf cover $\{S_i\to S\}$, every descent datum for
objects is effective.
\end{enumerate}
\end{Def}
This is Definition~8.5.1, \BGtag{02ZI}, with its fibration condition
already established. We spell out the terminology in the second clause.

\subsection{First picture: an ordinary open cover}
For the moment let the $S_i$ be open subsets covering $S$. Suppose a
double cover is given over each $S_i$. On $S_i\cap S_j$, an
isomorphism says which sheet in the first description corresponds to
which sheet in the second. On a triple intersection, changing from the
first description to the third directly must have the same result as
changing first to the second and then to the third. This is the
\emph{cocycle condition}: the identifications must not depend on the
route taken.

To \emph{descend} this information means to construct a single double
cover over $S$ with these local descriptions. The datum is called
\emph{effective} when such a global cover actually exists. For general
fppf covers, the same questions make sense, but $S_i\times_S S_j$
takes the place of the intersection. Its two projections keep track
of which local description is being used. These projections must be
retained because the covering maps need not be inclusions.

\begin{Ex}[Two routes through the sheet labels]
Suppose the local double covers have been labeled on the open sets
$S_i$. On a triple overlap $V=S_i\cap S_j\cap S_k$, each
description is $V\amalg V$; call the three labeled descriptions
$L_i,L_j,L_k$. A transition from labels $i$ to labels $j$ adds a
locally constant value $\varepsilon_{ij}$ in $\bZ/2\bZ$.
The two proposed routes from $L_i$ to $L_k$ are
\[
\begin{tikzcd}[column sep=small,row sep=large]
 L_i \arrow[rr,"+\varepsilon_{ik}"]
     \arrow[dr,"+\varepsilon_{ij}"']&&L_k\\
 &L_j \arrow[ur,"+\varepsilon_{jk}"']&
\end{tikzcd}
\]
Starting with label $a$, the lower route gives
$a+\varepsilon_{ij}+\varepsilon_{jk}$, and the direct route gives
$a+\varepsilon_{ik}$. The diagram commutes precisely when
\[
 \varepsilon_{ij}+\varepsilon_{jk}=\varepsilon_{ik}
 \quad\text{on }V.
\]
For instance, if the first two transitions swap sheets, the direct
transition must preserve them. Choosing all three transitions to be
swaps on a nonempty $V$ is inconsistent: the lower route sends $0$
to $0$, but the direct route sends $0$ to $1$. Each proposed map
individually is a valid isomorphism; it is their joint compatibility
that fails. This is the extra information recorded by a cocycle.
\end{Ex}

\subsection{The general datum}
Put $S_{ij}=S_i\times_S S_j$ with projections $r_1,r_2$, and put
$S_{ijk}=S_i\times_S S_j\times_S S_k$. On triple overlaps let $q_{ab}$
select factors $a,b$. A \emph{descent datum} consists of $x_i\in\cF_{S_i}$
and isomorphisms
\[
 \varphi_{ij}:r_1^*x_i\xrightarrow{\sim}r_2^*x_j
\]
such that
\begin{equation}\label{eq:cocycle}
 q_{23}^*\varphi_{jk}\circ q_{12}^*\varphi_{ij}
 =q_{13}^*\varphi_{ik}.
\end{equation}
All arrows in this equation are over $S_{ijk}$. This is
Definition~8.3.1, \BGtag{026B}, with projection indices starting at $1$.

The datum is \emph{effective} if there is $x\in\cF_S$ with
$\psi_i:x|_{S_i}\simeq x_i$ for which
\begin{equation}\label{eq:recovery}
 \varphi_{ij}=(r_2^*\psi_j)\circ(r_1^*\psi_i)^{-1}.
\end{equation}
The two pullbacks of $x$ in the middle are identified canonically.

\subsection{Self-overlaps require care}\label{subsec:selfoverlap}
An fppf cover need not consist of open immersions. Thus $S_{ii}$ need
not equal $S_i$, and the two objects $r_1^*x_i,r_2^*x_i$ on $S_{ii}$
are different pullbacks. The correct normalization is
\begin{equation}\label{eq:diagonal}
 \Delta_i^*\varphi_{ii}=\id_{x_i},
\end{equation}
where $\Delta_i:S_i\to S_i\times_S S_i$. If
$s_{ij}:S_{ij}\to S_{ji}$ exchanges factors, then
\[
 s_{ij}^*\varphi_{ji}=\varphi_{ij}^{-1}.
\]
To prove \eqref{eq:diagonal}, pull the $(i,i,i)$ cocycle back along the
small diagonal. The resulting invertible endomorphism $a$ satisfies
$a^2=a$, hence $a=\id$. Pulling the $(i,j,i)$ cocycle back along
$(u,v)\mapsto(u,v,u)$ then proves the inverse formula.
These are Remarks~8.3.2, \BGtag{026C}. The assertion
``$\varphi_{ii}=\id$ on all of $S_{ii}$'' is in general incorrect.

\subsection{The categorical meaning of descent}
Morphisms between descent data are families $u_i:x_i\to y_i$ such that
\[
 \eta_{ij}\circ r_1^*u_i=r_2^*u_j\circ\varphi_{ij},
\]
where $\eta_{ij}$ is the transition system for $y_i$. They form the
descent groupoid $\operatorname{Desc}_{\cF}(\{S_i\to S\})$.
Restriction is a functor
\[
 \cF_S\longrightarrow\operatorname{Desc}_{\cF}(\{S_i\to S\}).
\]
The Isom-sheaf axiom says it is fully faithful: compatible local arrows
between two existing global objects come from unique global arrows.
Effective descent says it is essentially surjective: each local object
system has a global realization. Both are required for an equivalence.

\BGcheckpoint{\begin{itemize}
\item Defined: the Isom-sheaf axiom, descent datum, cocycle, and
effective descent.
\item Check: state separately what must glue in each stack axiom.
\item Still to prove: both gluing assertions for $BG$.
\end{itemize}}

\begin{Ej}[A fibered groupoid with insufficient arrow descent]\label{ex:cfg}
On the open-set site of a space consisting of two discrete points, set
$A(U)=\bZ/2\bZ$ for nonempty $U$ and $A(\varnothing)=\{0\}$, with
the evident restrictions. Form the one-object groupoid with automorphism
group $A(U)$ over every open $U$. Explain why its associated total
category is fibered in groupoids but its automorphisms do not form a sheaf.
\end{Ej}

\section{Isomorphisms form a sheaf}\label{sec:isomsheaf}

\subsection{The presheaf and its exact restriction laws}
The objects $P$ and $Q$ are already defined over all of $S$. What will
vary locally is an \emph{identification between them}. Over a
$B$-scheme $T\to S$, such an identification is an equivariant
isomorphism $P_T\to Q_T$. Thus the elements of the set we are about
to define are maps, not points of either cover.

Restriction has a concrete meaning. Given $u:T'\to T$, pull back the
map $\varphi:P_T\to Q_T$ along $u$. On the corresponding fiber
products it sends $(p,t')$ to $(\varphi(p),t')$: the identification of
the covers is unchanged, but the base has been changed to $T'$.
The comparison maps below only identify these iterated fiber products
with the previously chosen notation $P_{T'}$ and $Q_{T'}$.

Fix $G_S$-torsors $P,Q$. Define
\[
\begin{aligned}
 \cI(T)&=\underline{\Isom}_G(P,Q)(T)\\
       &=\Isom_{G_T}(P_T,Q_T).
\end{aligned}
\]
For $u:T'\to T$ over $S$, let
\[
 c_{P,u}:P_{T'}\xrightarrow{\sim}u^*(P_T)
\]
be the canonical comparison, and similarly for $Q$. Define the
restriction function by
\[
 \rho_u(\varphi)=c_{Q,u}^{-1}\circ u^*\varphi\circ c_{P,u}.
\]
It is a map between the chosen objects $P_{T'}$ and $Q_{T'}$.
Identity comparisons give $\rho_{\id}=\id$. If $v:T''\to T'$, coherence
of the fiber-product comparisons gives
$\rho_v\rho_u=\rho_{uv}$ as an \emph{equality of functions}. Thus
$\cI$ is an actual set-valued presheaf, not a presheaf only up to
isomorphism. This construction is formalized in Lemma~8.2.1 and
Definition~8.2.2, \BGtag{026A} and \BGtag{02ZB}.

\subsection{The external gluing theorem}
\begin{Th}[Descent for scheme morphisms]\label{th:homdescent}
Let $X,Y$ be schemes over a scheme $S$, and let $\{S_i\to S\}$ be an
fpqc cover. Compatible $S_i$-maps $X_{S_i}\to Y_{S_i}$ glue to a unique
$S$-map $X\to Y$. In particular, equality of maps can be checked after
the cover. No affineness assumption on $X$ or $Y$ is required.
\end{Th}
This is Lemma~35.35.11 and its Hom-sheaf formulation in
Remark~35.35.12, \BGtag{02W0} and \BGtag{040L}. One underlying reason
is that representable functors are fpqc sheaves: $X_{S_i}\to X$ is a
cover, and compatible maps from these schemes to $Y$ glue by the sheaf
axiom for $h_Y$. Compatibility with the maps to $S$ is checked by the
same uniqueness statement.

\subsection{Apply the theorem to torsor isomorphisms}
The available data are an fppf cover $\{S_i\to S\}$ and equivariant
isomorphisms $\varphi_i:P_{S_i}\simeq Q_{S_i}$ agreeing on overlaps.
We must glue them and verify equivariance, invertibility, and uniqueness.

Theorem~\ref{th:homdescent} applies because $P,Q$ are schemes and the
cover is fpqc. It produces a unique underlying $S$-morphism
$\varphi:P\to Q$. To check equivariance, compare
\[
 \varphi a_P,\qquad a_Q(\id_G\times\varphi)
\]
as maps $G_S\times_S P\to Q$. Their base changes agree because every
$\varphi_i$ is equivariant. The equality part of the theorem makes
them equal globally.

The local inverses $\varphi_i^{-1}$ agree on overlaps and glue to
$\chi:Q\to P$. The maps $\chi\varphi$ and $\id_P$ agree on the cover,
as do $\varphi\chi$ and $\id_Q$. Hence they are equal globally and
$\varphi$ is an isomorphism. Its inverse is equivariant by the same
equation or by Lemma~\ref{lem:torsormap}. Finally, uniqueness of the
underlying map gives uniqueness among equivariant isomorphisms.

To prove the sheaf axiom on the \emph{whole} slice site, repeat this
argument over every $T\to S$ and every fppf cover of $T$, with $P_T,Q_T$
in place of $P,Q$.

\BGcheckpoint{\begin{itemize}
\item Constructed: the unique glued map and its inverse.
\item Verified: equivariance and both inverse equations.
\item Result: $\underline{\Isom}_G(P,Q)$ is an fppf sheaf, the first
stack axiom.
\end{itemize}}

For trivial double covers, a section of this sheaf is a locally constant
choice to preserve or exchange the two labels. Compatible choices on a
cover glue uniquely. For two nontrivial double covers, the same argument
glues their local maps commuting with the involutions; no global labeling
of either cover is required.

\begin{Ej}[The presheaf law]\label{ex:presheaf}
Using the comparisons $c_{P,u}$, verify $\rho_v\rho_u=\rho_{uv}$.
Explain why the restriction law is an equality even though iterated
fiber products are only canonically isomorphic.
\end{Ej}
\begin{Ej}[Gluing]\label{ex:glueiso}
Write the two maps whose equality expresses equivariance of a glued
$P\to Q$. Then prove invertibility by gluing local inverses. Identify
which clauses of Theorem~\ref{th:homdescent} are used for existence and
for uniqueness.
\end{Ej}

\section{Effective descent for torsors}\label{sec:effective}

\subsection{The affine descent input}
Here \emph{affine over $S$} is a property of the morphism to $S$;
it does not require the total space to be an affine scheme. For
example, $S\amalg S\to S$ is affine: above an affine open
$V=\Spec A$ it is
$V\amalg V=\Spec(A\times A)$. This remains true when
$S=\bP^1_k$, although $S\amalg S$ is not an affine scheme.
The affine descent theorem concerns this relative meaning of affine.

The remaining stack condition starts with local objects, so it must
construct a new global scheme. We use the following theorem as a black box.

\begin{Th}[Effective descent of affine morphisms]\label{th:affinedescent}
Let $\{S_i\to S\}$ be an fpqc cover. Suppose $X_i\to S_i$ are affine
morphisms with isomorphisms $r_1^*X_i\simeq r_2^*X_j$ satisfying the
cocycle condition. There exists an affine $S$-scheme $X$ and
isomorphisms $X_{S_i}\simeq X_i$ recovering those transitions.
Together with Theorem~\ref{th:homdescent}, this gives an equivalence
between affine $S$-schemes and their descent data.
\end{Th}
This is Lemma~35.37.1, \BGtag{0245}. The theorem depends on faithfully
flat descent of modules and algebras, not on torsor stackness. In affine
coordinates it descends an algebra with compatible identifications to
an algebra over the base ring; taking its spectrum constructs the
global affine scheme. A cocycle is needed so that these identifications
are mutually consistent.

We also use Lemma~35.23.20, \BGtag{02L5}: if a morphism of schemes
becomes affine over an fpqc cover of its target, it was affine already.
This assertion concerns a \emph{given} global morphism; it is distinct
from the effective-descent theorem, which constructs a global object.

\subsection{The datum and the underlying scheme}
Suppose given $G_{S_i}$-torsors $P_i\to S_i$ and equivariant
$\varphi_{ij}:r_1^*P_i\simeq r_2^*P_j$ satisfying
\eqref{eq:cocycle}. They prescribe how local torsors should be identified
on repeated base changes. We seek $(P,\psi_i)$ satisfying
\eqref{eq:recovery} and all torsor axioms.

Each $P_i\to S_i$ is affine: a trivializing fppf cover makes it
isomorphic to the affine morphism $G_{S_i}\to S_i$, and
\BGtag{02L5} applies. The outer cover is fpqc by \BGtag{022C}.
Thus every hypothesis of Theorem~\ref{th:affinedescent} is verified.
It gives an affine $P\to S$ and isomorphisms
\[
 \psi_i:P_{S_i}\xrightarrow{\sim}P_i
\]
recovering the specified transitions. At this stage $P$ is a scheme over
$S$; its action is still to be constructed.

\subsection{Descend the action and its equations}
Transport each action $a_i:G_{S_i}\times P_i\to P_i$ to
\[
 a_i'=\psi_i^{-1}a_i(\id_G\times\psi_i)
\]
on $(G_S\times_S P)_{S_i}\to P_{S_i}$. On $S_{ij}$, equivariance of
$\varphi_{ij}$ says
\[
 \varphi_{ij}a_i=a_j(\id_G\times\varphi_{ij}).
\]
Substitution of \eqref{eq:recovery} shows that $a_i'$ and $a_j'$ agree.
Theorem~\ref{th:homdescent}, applied with $X=G_S\times_S P$ and $Y=P$,
produces a unique map $a:G_S\times_S P\to P$.

Associativity compares the two maps from $G_S\times_S G_S\times_S P$
to $P$:
\[
 a(m\times\id_P),\qquad a(\id_G\times a).
\]
They agree after every $S_i$ because $a_i$ is an action. Thus they agree
globally. The same equality test proves
$a(e\times\id_P)=\id_P$. We have obtained a group action, and the
$\psi_i$ are equivariant by the way it was constructed.

\subsection{Local triviality and the torsor identity}
For each $i$, choose an fppf cover $\{U_{ia}\to S_i\}$ trivializing
$P_i$. Composition of fppf coverings makes
$\{U_{ia}\to S\}_{i,a}$ a cover. On it,
\[
 P_{U_{ia}}\simeq(P_i)_{U_{ia}}\simeq G_{U_{ia}}
\]
equivariantly. Consequently $P\to S$ is locally trivial in the chosen
topology. Proposition~\ref{prop:torsorcriterion} also gives its
surjectivity, flatness, and local finite presentation.

For an explicit check of the torsor identity, form
\[
 \Theta_P:G_S\times_S P\to P\times_S P.
\]
Under $\psi_i$ its base change is $\Theta_{P_i}$, hence an isomorphism.
The isomorphism-locality theorem \BGtag{02L4} therefore makes
$\Theta_P$ an isomorphism. This verifies the requested identity globally,
in addition to local triviality.

\subsection{Recovery and compatible uniqueness}
The maps $\psi_i$ already satisfy \eqref{eq:recovery}, by the affine
descent construction, and they are now maps of torsors. Suppose
$(P',\psi_i')$ is another realization. Define
\[
 \lambda_i=(\psi_i')^{-1}\psi_i:P_{S_i}\simeq P'_{S_i}.
\]
The two recovery equations imply agreement of $\lambda_i$ and
$\lambda_j$ on overlaps. Section~\ref{sec:isomsheaf} gives a unique
equivariant isomorphism $\lambda:P\simeq P'$ such that
\[
 \psi_i'\lambda_{S_i}=\psi_i
\]
for every $i$. The compatibility equations, not a claim that torsors have
no automorphisms, are what force uniqueness.

\subsection{The sheaf interpretation}
There is also a useful way to see the gluing before testing
representability. On a test $T\to S$, put $T_i=T\times_S S_i$ and
consider tuples
\[
 (p_i)\in\prod_i P_i(T_i)
\]
such that $\varphi_{ij}(p_i)=p_j$ on $T\times_S S_{ij}$.
These sets, with restriction of tuples, form a sheaf $\cP$.
Indeed, compatible tuples on a cover of $T$ glue componentwise by the
sheaf property of each $P_i$, and their compatibility equations can be
checked on that cover. Over $S_j$, restriction to the diagonal of the
$j$-th component gives $\cP|_{S_j}\simeq h_{P_j}$; the inverse
transports that section to the other components by $\varphi_{ji}$.
The cocycle proves compatibility and that these maps are inverse.
The action is componentwise.

This argument constructs a sheaf torsor. Affine descent is the separate
representability step, and identifies $\cP$ with $h_P$. Thus even this
version uses no circular appeal to the stack of torsors.

\BGcheckpoint{\begin{itemize}
\item Constructed: the global scheme, action, and local identifications.
\item Verified: action laws, local triviality, the torsor identity, and
recovery of transitions.
\item Result: effective descent, with uniqueness compatible with the
local identifications.
\end{itemize}}

For trivial double covers the transition maps are
$x\mapsto x+\varepsilon_{ij}$; their cocycle says the parity of sheet
swaps is independent of the route through a triple overlap. Affine
descent constructs the double cover, and equivariance glues its
involution. This remains valid for a singleton fppf cover with a
nontrivial self-overlap.

\begin{Ej}[A descent datum is more than local objects]\label{ex:datum}
Write a torsor descent datum using $r_1,r_2$ and $q_{12},q_{23},q_{13}$.
For trivial double covers interpret the cocycle as an equation for
$\varepsilon_{ij}$. On a triple overlap, why is the assignment
$\varepsilon_{ij}=\varepsilon_{jk}=1$, $\varepsilon_{ik}=1$ inconsistent?
\end{Ej}
\begin{Ej}[The self-overlap error]\label{ex:self}
Let $U=P\to S$ be a double-cover torsor and use its tautological
section to trivialize $P_U$. On $U\times_S U$, define
$\varepsilon(u,v)$ by $v=\varepsilon(u,v)\cdot u$.
Show that this is its transition function, that it vanishes on the
diagonal, and that it is $1$ on the graph of the deck involution.
Why does this refute $\varphi_{11}=\id$ on the entire self-overlap?
\end{Ej}
\begin{Ej}[Descend a structure map]\label{ex:action}
Assume affine descent has produced $P$ and $\psi_i$. Prove in detail
that the transported actions $a_i'$ agree, and list the domains of the
two morphisms used to check associativity. Is checking the action only
on geometric points sufficient?
\end{Ej}
\begin{Ej}[Uniqueness of a realization]\label{ex:realization}
For two realizations $(P,\psi_i)$ and $(P',\psi_i')$, prove that
$(\psi_i')^{-1}\psi_i$ agrees on overlaps. State exactly the uniqueness
assertion that follows, and explain why omitting the $\psi_i$ changes it.
\end{Ej}

\section{The stack theorem}\label{sec:theorem}

We now apply Definition~\ref{def:stack} to the structure we have built.
The required data and their verifications are all available:
\begin{enumerate}[label=(\arabic*)]
\item Section~\ref{sec:total} established an ordinary category $BG$.
\item Section~\ref{sec:projection} established the functor to $\Sch/B$.
\item Sections~\ref{sec:pullback}--\ref{sec:cartesian} constructed
cartesian lifts of all base maps.
\item Section~\ref{sec:groupoids} proved invertibility in every fiber.
\item Section~\ref{sec:isomsheaf} proved the Isom-sheaf condition on
every slice site.
\item Section~\ref{sec:effective} proved effective object descent.
\end{enumerate}
These prove Theorem~\ref{th:target}.

The functor from the naive prestack to this torsor category is fully
faithful on fibers: its arrows are exactly the translations described
in Section~\ref{sec:quotient}. Every torsor is locally in its image.
Since the target is now proved to be a stack, the stackification
comparison of \BGtag{0CQJ} identifies it with $[B/G]$.

\BGcheckpoint{\begin{itemize}
\item Verified: a category over $\Sch/B$, cartesian lifts, and groupoid
fibers.
\item Verified: both descent axioms.
\item Result: $BG=[\operatorname{pt}/G]$ is a stack for the fppf
covering rule.
\end{itemize}}

For the constant group of order two, the theorem says that equivariant
isomorphisms of double covers glue uniquely and that every compatible
system of local double covers has a global realization.

\section{A detailed double-cover example}\label{sec:double}

\subsection{The groupoid of double covers}
Let $G=(\bZ/2\bZ)_B$. A torsor is finite \emph{\'etale} of degree $2$:
it has this property on a trivializing cover, and finite locally free
morphisms of fixed degree and \emph{\'etale} morphisms descend
fpqc-locally (Lemmas~35.23.32 and 35.23.31, \BGtag{02VO} and
\BGtag{02VN}).
Conversely, a finite \emph{\'etale} degree-two cover is
\emph{\'etale}-locally two copies of its base
(Lemma~41.18.3, \BGtag{04HN}). The unique permutation
exchanging the two points commutes with every permutation of two points,
so its local involutions glue, independently of labeling, to a canonical
action of $\bZ/2\bZ$. This is a torsor action.

For the trivial cover over $S$, an automorphism is
$x\mapsto x+\varepsilon$ with $\varepsilon\in G(S)$. If $S$ is
nonempty and connected, there are exactly two choices. On a disconnected
base the choice can differ on each open-and-closed component; it is
incorrect to count exactly two global automorphisms for every base.

\subsection{A nonsplit cover and its splitting cover}
Take a field $k$ with $\Char k\ne2$, let
\[
\begin{aligned}
 S&=\Spec k[t,t^{-1}],\\
 P&=\Spec k[u,u^{-1}],
\end{aligned}
\]
and use $t=u^2$ for the structural map. The coordinate algebra is free
of rank two over $k[t,t^{-1}]$, with basis $1,u$, and the derivative
$2u$ is a unit. The standard \emph{\'etale} algebra criterion says that
a monic polynomial with invertible derivative in its quotient algebra
defines an \emph{\'etale} algebra (Definition~10.144.1 and
Lemma~10.144.2(1), \BGtag{00UB} and \BGtag{00UC}). Thus our map is
finite \emph{\'etale}. The action is
$u\mapsto-u$.

The torsor identity can be checked without counting points: the algebra
of $P\times_S P$ is
\[
 k[u^{\pm1},v^{\pm1}]/(u^2-v^2).
\]
The ideals $(u-v)$ and $(u+v)$ are comaximal because $2v$ is invertible.
Thus this scheme is the disjoint union of the diagonal and the graph of
the involution, exactly the two components of $G_S\times_S P$.

A section of $P\to S$ would give $w\in k[t,t^{-1}]$ with $w^2=t$.
Every unit has the form $ct^n$, so this would force $2n=1$.
There is no section, and hence the torsor is nontrivial by
Exercise~\ref{ex:section}. Nevertheless it splits on the single fppf
cover $P\to S$. Its self-overlap records whether $v=u$ or $v=-u$.
The transition $\varepsilon(u,v)$ equals $0$ or $1$ respectively, and
on a triple overlap
\[
 \varepsilon(u,v)+\varepsilon(v,w)=\varepsilon(u,w).
\]
The equality expresses the composition of signs. Descent of the trivial
torsor over $P$, with this transition, reconstructs the original
nonsplit torsor over $S$.

\subsection{Characteristic two}
Over a field $k$ of characteristic $2$, the morphism
\[
 \Spec k[x]\longrightarrow\Spec k[t],
 \qquad t=x^2-x,
\]
is free of degree two, with invertible derivative $-1$.
Its involution is $x\mapsto x+1$. On the self-product the difference
$y-x$ satisfies
\[
 (y-x)^2-(y-x)=0,
\]
which splits into the comaximal factors $y-x$ and $y-x-1$.
This verifies the constant-group torsor identity in characteristic $2$.
A polynomial section would solve $r(t)^2-r(t)=t$; a nonconstant $r$
has even degree on the left and a constant $r$ gives a constant. Thus
the torsor is nontrivial.

By comparison, over $k[t,t^{-1}]$ the map $t=u^2$ in characteristic $2$
is a $\mu_2$-torsor: it is finite faithfully flat and the ratio of two
lifts is an element of $\mu_2$. It is not \emph{\'etale}, and the
formula $u\mapsto-u$ becomes the identity, so it does not define a
constant $\bZ/2\bZ$-torsor action.

\begin{Ej}[Automorphisms on a disconnected base]\label{ex:automorphisms}
Let $S=S_1\amalg S_2$, with both $S_i$ nonempty and connected. Determine
the automorphism group of the trivial constant-group double cover over
$S$. Explain its sheaf interpretation.
\end{Ej}
\begin{Ej}[Artin--Schreier check]\label{ex:AS}
Verify the degree, \emph{\'etale} property, involution, torsor identity,
and nontriviality of $t=x^2-x$ in characteristic $2$. Locate precisely
where the analogous formula $u\mapsto-u$ fails for the constant group.
\end{Ej}

\section{Comparison with families of curves}\label{sec:curves}

\subsection{The actual moduli problem}
The Stacks Project studies curves in Section~99.15, \BGtag{0D4Y}.
Its category $\Curves$ (Situation~99.15.1, \BGtag{0D4Z}) has objects
$X\to S$ with $S$ a scheme, $X$ an algebraic space, and the morphism
flat, proper, finitely presented, of relative dimension at most one.
No smoothness, connectedness, reducedness, or fixed genus is imposed.
Relative dimension is in the sense of Definition~67.33.2,
\BGtag{06LR}. Here is the distinction from the dimension of the total
space: a family over a curve may have a two-dimensional total space,
while each fiber is only one-dimensional. The relative-dimension
condition concerns those fibers. More precisely, at a point of the
family one measures the dimension there in the fiber over its image
in the base; the bound must hold at every point. Thus ``at most one''
also allows zero-dimensional fibers, rather than requiring every fiber
to be a smooth, connected curve. We may restrict the base schemes to
$B$-schemes throughout.

An arrow is a \emph{cartesian} square
\[
\begin{tikzcd}[column sep=large,row sep=large]
 X' \arrow[r] \arrow[d] & X \arrow[d]\\
 S' \arrow[r,"f"'] & S
\end{tikzcd}
\]
If all commutative squares were allowed, arrows in a fixed fiber could
include noninvertible endomorphisms; the fibers would not be groupoids.

\subsection{The formal categorical correspondence}
\noindent
\begin{tabularx}{\linewidth}{@{}>{\raggedright\arraybackslash}X>{\raggedright\arraybackslash}X@{}}
\toprule
$BG$ & $\Curves$\\
\midrule
Torsor $P\to S$ & Family $X\to S$\\
Equivariant arrow over $f$ & Cartesian square over $f$\\
$P\simeq f^*Q$ & $X'\simeq f^*X$\\
Pullback torsor & Base-changed family\\
Equivariant isomorphism & Isomorphism of families\\
Affine descent & Algebraic-space descent\\
\bottomrule
\end{tabularx}
\par\smallskip

Identities and composites are defined componentwise; cartesian squares
paste, so the result is a category. Forgetting $X$ gives a functor to
schemes. Base change gives the cartesian lift, by exactly the same
factorization proof as Section~\ref{sec:cartesian}. Flatness, properness,
finite presentation, and the relative-dimension bound are stable under
arbitrary base change. For the first three properties see
Lemmas~67.30.4, 67.40.3, and 67.28.3, respectively:
\BGtag{03MO}, \BGtag{04WP}, and \BGtag{03XR}. For the dimension bound
use Lemma~67.34.3(1), \BGtag{04NS}. Finally a cartesian square over
$\id_S$ identifies $X'$
with $S\times_S X=X$, so each fiber is a groupoid.

The base-change and locality facts used here are stated below. Thus the
formal categorical argument is identical, while closure of the class
of allowed objects is a geometric theorem.

\subsection{Isomorphism descent for curves}
For global $X,Y\to S$, consider
\[
 T\longmapsto\Isom_T(X_T,Y_T).
\]
The relevant black box is Lemma~95.7.2, \BGtag{04UA}: compatible local
morphisms between two algebraic spaces over a scheme glue uniquely on
an fppf covering. It applies to these $X,Y$ even if they are not schemes.
Gluing inverses and checking the two identity composites on the cover
proves the Isom-sheaf condition. No representability theorem for this
Isom sheaf is needed for stackness.

\subsection{Effective descent and the geometric conditions}
The required object-descent theorem is Lemma~80.11.3(2),
\BGtag{0ADV}: descent data for algebraic spaces $X_i\to S_i$ of finite
type are effective for any fppf covering family $\{S_i\to S\}$.
It yields an algebraic space $X\to S$ and compatible
$X_{S_i}\simeq X_i$. Finite presentation implies finite type, so it
applies to our families. We use part~(2) explicitly; the unrestricted
part~(1) has a countability qualification in the Stacks Project's size
conventions.

It remains to verify that $X\to S$ is an allowed object. For algebraic
spaces the following properties are fpqc-local on the target, so they
are detected on our cover:
\begin{itemize}
\item finite presentation:\newline Lemma~74.11.12, \BGtag{041V};
\item flatness:\newline Lemma~74.11.13, \BGtag{041W};
\item properness:\newline Lemma~74.11.19, \BGtag{0422}.
\end{itemize}
Thus $X\to S$ has these three properties. In particular it is locally
of finite type. Lemma~67.34.3(1), \BGtag{04NS}, says that its relative
dimension at a point equals that of any base change at a point above it.
Every point of $X$ lifts to some $X_{S_i}$ because the base-changed
cover is jointly surjective. The local bound at most one therefore
holds globally. Preservation under base change alone would not supply
this converse; joint surjectivity is also used.

The recovery maps are those of the descent theorem, and two realizations
are uniquely isomorphic compatibly by morphism descent. This proves the
two missing stack axioms for $\Curves$ using explicit geometric inputs.
It is the argument packaged in Lemma~99.13.3, \BGtag{0D1G}, and
Lemma~99.15.3, \BGtag{0D52}; these packaged results are checks on our
argument, not assumptions that $\Curves$ is already a stack.

\subsection{Why algebraic spaces occur}
Schemes have fully faithful fpqc descent, but arbitrary scheme descent
data need not be effective as schemes. Lemma~110.68.1, \BGtag{08KF},
gives even projective schemes with \emph{\'etale} descent data that fail
to descend in schemes. This is a general example, not specifically a
curve-family example. For curves themselves, Lemma~110.69.1,
\BGtag{0D5E}, gives a family over $\bA^1_k$ with total space not a scheme.
Allowing algebraic spaces is therefore substantive.

For $BG$, affineness made both the local objects and their descended
object schemes. For curves, effective descent uses the larger category
of algebraic spaces. The Isom axiom only glues arrows between existing
objects; the object axiom must construct geometry and check it remains
in the class under consideration. This explains why the latter is
usually the harder assertion.

\begin{Ej}[Transfer the proof to curves]\label{ex:curves}
For each step from category to stack, identify the corresponding step
for $\Curves$. Name the extra geometric inputs. Detect the errors in:
``Every commutative square of curve families is an arrow'' and
``Morphisms of schemes descend, so local families always glue to a
scheme.''
\end{Ej}

\section{Summary of the logical progression}\label{sec:summary}

The following chain records the cumulative construction; each downward
step includes a new verification, rather than claiming an implication
from one isolated property to the next.
\[
\boxed{\begin{gathered}
 BG\text{ is a category}\\ \Downarrow\\
 BG\text{ is over }\Sch/B\\ \Downarrow\\
 BG\text{ is fibered over }\Sch/B\\ \Downarrow\\
 BG\text{ is fibered in groupoids}\\ \Downarrow\\
 \underline{\Isom}_G\text{ is a sheaf}\\ \Downarrow\\
 \text{torsor descent is effective}\\ \Downarrow\\
 BG\text{ is a stack}.
\end{gathered}}
\]
Here is the inventory of black boxes, separated from what the proofs
constructed explicitly.
\begin{enumerate}[label=(\arabic*)]
\item Fiber products of schemes and their universal property; base
change and composition of fppf covers, \BGtag{021O}; and the inclusion
of fppf covers among fpqc covers, \BGtag{022C}.
\item Subcanonicity and full faithfulness of scheme descent:
\BGtag{023Q}, \BGtag{02W0}, and \BGtag{040L}.
\item Locality of isomorphisms, affineness, flatness, and local finite
presentation: \BGtag{02L4}, \BGtag{02L5}, \BGtag{02L2}, \BGtag{02KY}.
\item Effective descent of affine morphisms, \BGtag{0245}. This is the
substantial existence theorem for the torsor proof. The action and its
equations were then descended explicitly.
\item For the interpretation of notation, the quotient-stack
definition and torsor comparison, \BGtag{044Q} and \BGtag{0CQJ}.
These are not used to prove torsor descent.
\item For the curve comparison, morphism descent \BGtag{04UA},
effective descent \BGtag{0ADV}(2), and locality of the defining
properties \BGtag{041V}, \BGtag{041W}, \BGtag{0422}, \BGtag{04NS}.
Their base-change stability uses \BGtag{03MO}, \BGtag{04WP},
\BGtag{03XR}, and \BGtag{04NS} separately from locality.
\end{enumerate}
The double-cover computations use the elementary finite
\emph{\'etale} algebra criterion (a monic polynomial with invertible
derivative) and finite \emph{\'etale} local splitting. The
algebra criterion is \BGtag{00UB}--\BGtag{00UC}; the splitting result
is \BGtag{04HN}. Descent of finite locally free morphisms of fixed
degree and of \emph{\'etale} morphisms uses \BGtag{02VO} and
\BGtag{02VN}. Smoothness locality and local sections, used only to
compare topologies, are \BGtag{02VL} and \BGtag{055U}. The constant
group and $\mu_2$ were distinguished by their coordinate algebras, not
by a characteristic-independent identification.

\section{Exercises: synthesis and review}\label{sec:exercises}

Exercises~\ref{ex:fiber}--\ref{ex:curves} appear immediately after the
relevant constructions. The next two combine several stages; the final
six return to the foundations and can be attempted on a first reading.
All exercises have a numbered solution in Section~\ref{sec:solutions}.

\begin{Ej}[Change of trivializations]\label{ex:change}
Let $\tau_i:P_{S_i}\simeq G_{S_i}$ be left-equivariant trivializations,
and write $\tau_j\tau_i^{-1}=R_{h_{ij}}$. Prove
\[
 h_{ik}=h_{ij}h_{jk}
\]
on triple overlaps. If $\tau_i'=R_{b_i}\tau_i$, compute $h_{ij}'$ and
prove the new transitions satisfy the same cocycle. Interpret the
answer for $G=\bZ/2\bZ$.
\end{Ej}
\begin{Ej}[Two failures of forgetting data]\label{ex:forget}
Take $B=\Spec\bR$ and $G=(\bZ/2\bZ)_B$. Compare the split torsor
with $\Spec\bC\to\Spec\bR$, where the involution is complex
conjugation. Prove that they become isomorphic after the fppf cover
$\Spec\bC\to\Spec\bR$ but are not isomorphic over $\bR$. Deduce:
\begin{enumerate}[label=(\alph*)]
\item the prestack of only trivial torsors can satisfy the Isom-sheaf
condition while failing effective object descent;
\item taking isomorphism classes of all torsors gives a presheaf that
need not be a sheaf, although the torsor category is a stack.
\end{enumerate}
\end{Ej}

\subsection{Foundational review}
\begin{Ej}[Disjoint union versus product]\label{ex:coproducts}
Take $B=\Spec k$ and $S=\bA^1_k$. Give the coordinate rings of
$S\amalg S$, $S\times_B S$, and $S\times_S S$. Describe the ring
map defining the fold $S\amalg S\to S$. Explain why the fold is an
arrow of $\Sch/B$ and of $\Sch/S$, and identify its two sections.
\end{Ej}

\begin{Ej}[A map into a fiber product]\label{ex:universalmap}
Take $X=T=\bA^1_k$ over $S=\Spec k$, and let $W=\Spec k$ map to
the origin of each. Draw the induced map $W\to T\times_S X$ and
explain why it is not an isomorphism. What extra property of $W$
would make the comparison an isomorphism? Separately, pull the family
$x^2=t$ back by $t=1$ and compute the result when $\Char k\ne2$.
\end{Ej}

\begin{Ej}[Test points and Yoneda in coordinates]\label{ex:testpoints}
Let $T=\Spec k[\epsilon]/(\epsilon^2)$. Show that sending the
coordinate of $\bA^1_k$ to $0$ and to $\epsilon$ gives different
$T$-points with the same image on topological points. For every
$k$-algebra $R$, the rule $r\mapsto r^2$ is a function on the
$R$-points of $\bA^1_k$. Check naturality and recover the scheme
morphism by evaluating on the identity point, as in the Yoneda proof.
\end{Ej}

\begin{Ej}[Three meanings of a group symbol]\label{ex:groupobjects}
For $H=\bZ/2\bZ$ and $B=\bA^1_k$, distinguish $H$, $H_B$, and
$H_B(T)$. Describe the unit, multiplication, and inversion of $H_B$
on its labeled components. Compute $H_B(T)$ for $T=B$ and for
$T=B\amalg B$, and describe the pullback of $H_B$ to any $S\to B$.
\end{Ej}

\begin{Ej}[Finite presentation is a relative algebra condition]\label{ex:fpflat}
Explain why $A[z]/(z^2-z)$ is a finitely presented $A$-algebra and
flat as an $A$-module. Why is $A[x]$ finitely presented as an algebra
even though, for $A\ne0$, it is not finitely generated as an
$A$-module? Finally show that $\bZ/2\bZ$ is finitely presented as a
$\bZ$-algebra but is not flat over $\bZ$.
\end{Ej}

\begin{Ej}[The covering rule in use]\label{ex:coveroperations}
Given an fppf cover $\{S_i\to S\}$, a map $T\to S$, and covers
$\{U_{ij}\to S_i\}_j$, write the pulled-back and composite covering
families. Identify the relevant parts of \BGtag{021O} and explain
why \BGtag{021N} is not the reference for these assertions.
For the singleton fold cover $S\amalg S\to S$, compute its pullback
to $T$, and then cover each of its two sheets by two copies again.
Does every arrow in $(\Sch/B)_{\mathrm{fppf}}$ have to be flat?
\end{Ej}

\end{multicols}
\clearpage
\begin{multicols}{2}
\raggedcolumns
\section{Hints and solutions}\label{sec:solutions}

The exercise statements occur earlier so that they can be attempted
independently. Each solution begins with a short hint and then gives
the complete argument.

\BGsolution{ex:fiber}
\hint{First select the objects; then decide which arrows to retain.}
Both categories have the object collection $\{x:p(x)=S\}$. For two
objects in this collection, fullness means keeping every arrow:
\[
 \Hom_{\cD}(x,y)=\Hom_{\cF}(x,y).
\]
The fiber instead keeps precisely the subset on which $p(u)=\id_S$.
Since the source and target of $p(u)$ are both $S$, a discarded arrow
can project to a nonidentity endomorphism of $S$.

When $p=\id_{\cC}$, the equation $p(x)=S$ forces $x=S$. Thus there
is one object in each category. Their arrow sets are
\[
\begin{aligned}
 \End_{\cD}(S)&=\End_{\cC}(S),\\
 \End_{\cF_S}(S)&=\{\id_S\}.
\end{aligned}
\]
In particular $\varepsilon$ is present in $\cD$ and absent from the
fiber: the identity \emph{functor} leaves $\varepsilon$ unchanged;
it does not turn it into the identity \emph{morphism}. For
$S=\bA^1_k$, the maps corresponding to $t\mapsto t^2$ and
$t\mapsto t$ differ as ring maps, hence as scheme maps. This remains
true over a finite field even if those maps happen to agree on some
set of rational points.

There is no contradiction and no extra triangle condition for an
arbitrary functor. Both categories are subcategories with the same
objects, but only the first is required to be full. The word
``over $S$'' for an object is weaker than the condition ``over
$\id_S$'' for an arrow.

\BGsolution{ex:empty}
\hint{Both sides of the proposed identity are empty.}
Both $G_S\times_S\varnothing$ and
$\varnothing\times_S\varnothing$ equal $\varnothing$, so the map
between them is an isomorphism. But the empty morphism is not surjective
over nonempty $S$. No fppf cover can turn it into $G$, which has a unit
section. The missing requirement is local nonemptiness, supplied here
by local triviality or by the surjectivity of an fppf covering.

\BGsolution{ex:section}
\hint{Evaluate $\Theta_P^{-1}$ on $(p,s(\pi_P(p)))$.}
If $s$ is a section, the map $g\mapsto g\cdot s$ is equivariant. The
torsor identity gives, for every test point $p$, a unique $g$ with
$p=g\cdot s$, naturally in the test scheme, and hence gives a morphism
inverse. A trivial torsor has its unit section. If $s'=h\cdot s$ for
$h\in G(S)$, the two trivializations $F_s(g)=g\cdot s$ obey
$F_{s'}=F_s\circ R_h$. Every other section has this form, by the same
identity.

\BGsolution{ex:mu}
\hint{Compute idempotents and nilpotents.}
The constant group has algebra $k\times k$ and two $k$-points.
In characteristic $2$, $\mu_2$ has algebra
$k[\epsilon]/(\epsilon^2)$ and only one $k$-point: any field element
whose square is zero is zero. It also has a nonzero nilpotent element,
so its scheme structure differs from the reduced algebra $k\times k$.
Both algebras have dimension two as $k$-vector spaces, but rank does not
determine an algebra or a group scheme.

\BGsolution{ex:opposite}
\hint{Apply the composites to a variable $x$.}
We have $(R_k\circ R_h)(x)=(xh)k=x(hk)$.
Usual groupoid composition sends $(g_1,g_2)$ to $g_2g_1$, but
$(R_{g_2}\circ R_{g_1})(x)=xg_1g_2$. This reverses the required order.
Replacing each parameter by its inverse gives
$g_1^{-1}g_2^{-1}=(g_2g_1)^{-1}$, exactly the required functorial law.

\BGsolution{ex:axioms}
\hint{Compute the two components separately.}
For $a=(f,\widetilde f)$, $b=(g,\widetilde g)$,
$c=(h,\widetilde h)$, both $(cb)a$ and $c(ba)$ equal
$(hgf,\widetilde h\widetilde g\widetilde f)$. Both left and right
composition with the appropriate identity leave each component fixed.
The composites remain equivariant by the calculation in
Section~\ref{sec:total}. Because the top components are actual scheme
maps, this is strict associativity of maps. If arrows were specified
instead by $P\simeq f^*Q$, composition would require the comparison
$f^*g^*R\simeq(gf)^*R$.

\BGsolution{ex:arrows}
\hint{Evaluate at the unit section of $G_S$.}
The composite of $e_S:S\to G_S$ with the given arrow is a map $S\to G_T$
over $f$, hence a section $h\in G(S)$. Equivariance forces the formula
in the exercise, and that formula defines an equivariant morphism for
every $h$. Under the second arrow, $xh(s)$ is multiplied on the right
by $k(f(s))$. Thus the composite over $gf$ has parameter
$h\cdot f^*k\in G(S)$, in this order.

\BGsolution{ex:projection}
\hint{Use the unit parameter in Exercise~\ref{ex:arrows}.}
Take $\widetilde f=\alpha_f:G_S\to G_S$. The square commutes, and
$\alpha_f$ intertwines the two regular actions over $f$. The pair
$(f,\alpha_f)$ is therefore a total arrow. For a nonempty base the
endomorphism $t\mapsto t^2$ of $\bA^1_B$ differs from the identity,
so this arrow is not vertical and is excluded from the fiber. If $B$
is empty the example is vacuous; choose a nonempty base as intended.

\BGsolution{ex:pullback}
\hint{Factor $u^2-v^2$ and use that $2v$ is invertible.}
The coordinate algebra is
\[
 k[v^{\pm1},u]/(u^2-v^2).
\]
Here $u$ is automatically invertible. The ideals $(u-v)$ and $(u+v)$
are comaximal, so the Chinese remainder theorem identifies the algebra
with $k[v^{\pm1}]\times k[v^{\pm1}]$. The two sections are $u=v$ and
$u=-v$. The involution $u\mapsto-u$ exchanges these sections, giving
the trivial double cover of $T$.

\BGsolution{ex:cartesian}
\hint{There is no freedom once the two projections are prescribed.}
The projections must be $\pr_P\lambda=\widetilde h$ and
$\pr_T\lambda=k\pi_R$. They are compatible because
$\pi_P\widetilde h=fk\pi_R$. The fiber product gives a unique
$\lambda$. Applying it to $g\cdot r$ changes the first component by
$\alpha_h(g)$ and leaves the second unchanged. The induced action on
$P\times_S T$ uses $\alpha_f\alpha_k(g)$, which is the same element.
Thus the factorization is equivariant. Another arrow over $k$ satisfying
the factorization equation has both the same projections, hence is equal.

\BGsolution{ex:unique}
\hint{A pullback includes its arrow to the original object.}
For a cartesian lift $u:x\to y$, uniqueness concerns vertical
automorphisms $a$ with $ua=u$. The cartesian property forces such $a$
to be $\id_x$. For a pulled-back double cover the deck involution changes
the projected point of the original cover, so generally $ua\ne u$.
It is a torsor automorphism but not an automorphism of the prescribed
cartesian lift. The latter includes the structure arrow $u$.

\BGsolution{ex:invertible}
\hint{Common refinements trivialize both source and target.}
If $U_i$ trivializes $P$ and $V_j$ trivializes $Q$, then
$U_i\times_S V_j$ jointly form an fppf cover of $S$. On this cover the
map is $R_h$ for its value at the unit, hence has inverse $R_{h^{-1}}$.
Isomorphism locality proves global invertibility and equivariance of
the inverse follows from $u(gp)=gu(p)$. Geometric points alone do not
detect scheme isomorphisms: for example, $\Spec k\to
\Spec k[\epsilon]/(\epsilon^2)$ is a bijection on geometric points
but not an isomorphism. Finally, the pullback of a trivial torsor along
a proper nonempty open immersion of bases gives a total arrow whose
top map is an open immersion, not an isomorphism.

\BGsolution{ex:cfg}
\hint{Choose different automorphisms on the two singleton opens.}
Restrictions of $A$ are group homomorphisms. The associated total
category has one object over each open, and pullback along inclusions
is given by the designated object and the identity group element; the
groupoid construction yields cartesian lifts. Every vertical arrow is
invertible. On the two singleton opens choose automorphisms $0$ and $1$.
They agree on the empty intersection since $A(\varnothing)$ is trivial.
But a global section of $A$ is one element that restricts to the same
element on both singleton opens. No section has the prescribed pair
$(0,1)$, so the automorphism presheaf is not a sheaf.

\BGsolution{ex:presheaf}
\hint{Expand both restrictions and cancel adjacent comparisons.}
The composite restriction is
\[
\begin{aligned}
 \rho_v\rho_u(\varphi)
 ={}&c_{Q,v}^{-1}\,v^*(c_{Q,u}^{-1})\\
 &\circ v^*u^*\varphi\circ v^*c_{P,u}\circ c_{P,v}.
\end{aligned}
\]
Coherence identifies the outer composites with the comparisons for
$uv$, and identifies the middle map with $(uv)^*\varphi$.
Thus the result equals $c_{Q,uv}^{-1}(uv)^*\varphi\,c_{P,uv}$.
Both sides are maps between the same chosen $P_{T''}$ and $Q_{T''}$,
so this is equality, not an unspecified isomorphism of arrows.

\BGsolution{ex:glueiso}
\hint{Equivariance and inverse laws are equalities of morphisms.}
The maps are $\varphi a_P$ and $a_Q(\id_G\times\varphi)$ from
$G_S\times_S P$ to $Q$. The theorem's existence clause glues the
underlying $\varphi_i$, and its uniqueness clause makes these two
maps equal since they are equal after the cover. Apply existence again
to the compatible $\varphi_i^{-1}$, obtaining $\chi$. Apply uniqueness
to $\chi\varphi$ and $\id_P$, then to $\varphi\chi$ and $\id_Q$.
This proves invertibility. A second global isomorphism with the same
restrictions has the same underlying map by uniqueness.

\BGsolution{ex:datum}
\hint{All three maps in the cocycle must be over one triple overlap.}
The datum is $P_i$ and
$\varphi_{ij}:r_1^*P_i\simeq r_2^*P_j$, with
$q_{23}^*\varphi_{jk}\,q_{12}^*\varphi_{ij}=q_{13}^*\varphi_{ik}$.
For the constant group the transition is addition of
$\varepsilon_{ij}$, so composition is addition of the two pulled-back
sections. The required equation is
$\varepsilon_{ij}+\varepsilon_{jk}=\varepsilon_{ik}$.
Two swaps give $0$, not $1$, so the proposed triple gives incompatible
identifications whenever that triple overlap is nonempty.

\BGsolution{ex:self}
\hint{Compare the two tautological sections on the self-overlap.}
Use coordinates on $P_U$ given by the map $(x,u)\mapsto x\cdot u$.
On $U\times_S U$, the coordinate $x$ in the first trivialization and
$y$ in the second describe the same point exactly when
$x\cdot u=y\cdot v$. With $v=\varepsilon(u,v)\cdot u$, this says
$y=x+\varepsilon(u,v)$ in $\bZ/2\bZ$. The torsor identity makes
$\varepsilon$ a well-defined section. It is $0$ when $v=u$ and $1$
when $v=\sigma(u)$. Simple transitivity gives the cocycle by composing
these comparisons. Thus the transition is identity along the diagonal
but is a swap on the other component. A single covering member still
has meaningful descent data on its self-overlap.

\BGsolution{ex:action}
\hint{Substitute $\varphi_{ij}=\psi_j\psi_i^{-1}$ into equivariance.}
The equation $\varphi_{ij}a_i=a_j(\id\times\varphi_{ij})$, after
precomposing with $\id\times\psi_i$ and postcomposing with
$\psi_j^{-1}$, becomes
$\psi_i^{-1}a_i(\id\times\psi_i)
=\psi_j^{-1}a_j(\id\times\psi_j)$. These are the transported actions,
so they glue. Both associativity maps have domain
$G_S\times_S G_S\times_S P$ and codomain $P$; the unit equation has
domain $P$. Their equality is checked as maps after the fppf cover.
Geometric-point checks alone miss nilpotents, as in
Exercise~\ref{ex:invertible}, so they are insufficient.

\BGsolution{ex:realization}
\hint{Both local systems have the same transition map.}
On an overlap, the recovery equations yield
$\psi_j=\varphi_{ij}\psi_i$ and
$\psi_j'=\varphi_{ij}\psi_i'$, with pullbacks understood. Hence
\[
 (\psi_j')^{-1}\psi_j
 =(\psi_i')^{-1}\varphi_{ij}^{-1}\varphi_{ij}\psi_i
 =(\psi_i')^{-1}\psi_i.
\]
The Isom sheaf gives a unique $\lambda:P\simeq P'$ with
$\psi_i'\lambda_{S_i}=\psi_i$. If the local identification maps are
forgotten, this equation is no longer required, and one can compose
with any torsor automorphism. A trivial double cover already has a
nonidentity example.

\BGsolution{ex:automorphisms}
\hint{Choose independently on the two connected components.}
Each component admits the identity or the swap, so the global group is
$(\bZ/2\bZ)^2$ with four elements. In general the automorphism sheaf
of the trivial torsor is the sheaf of locally constant
$\bZ/2\bZ$-valued functions. The value on a disjoint union is the
product of the values on its components, exactly as the sheaf axiom
requires.

\BGsolution{ex:AS}
\hint{Subtract the two roots in the self-product.}
As a $k[t]$-algebra, $k[x]=k[t][x]/(x^2-x-t)$ is free with basis $1,x$.
The derivative is $-1$, so the finite flat map is \emph{\'etale}.
Translation by $1$ preserves $x^2-x$ and has order two. In the
self-product the difference $y-x$ satisfies $d(d-1)=0$, and $(d)$ and
$(d-1)$ are comaximal. These two components give the torsor identity.
A section would yield $r^2-r=t$ in $k[t]$. For positive degree $n$, the
left side has degree $2n$; for degree zero it is constant. Neither equals
$t$. Finally, $-u=u$ in characteristic $2$, so negation does not
exchange two sheets and cannot supply the nonidentity constant-group
action.

\BGsolution{ex:curves}
\hint{Separate universal properties from closure under descent.}
Define objects and cartesian arrows, then check identity and pasting of
squares to obtain a category. Forgetting total spaces is a functor.
The fiber product supplies cartesian lifts, while base-change stability
of the four geometric conditions ensures that these are allowed
objects. A fiber arrow is an isomorphism because its square is cartesian
over an identity. The morphism-descent theorem for algebraic spaces,
\BGtag{04UA}, glues local isomorphisms and their inverses. Effective
descent \BGtag{0ADV}(2) constructs the global algebraic space; the
locality theorems for flatness, properness, finite presentation, and
relative dimension verify its membership in $\Curves$.
An arbitrary commutative square need not be cartesian, so the first
proposed assertion has too many arrows. Full faithfulness for schemes
does not give essential surjectivity, and \BGtag{08KF} disproves the
second assertion. Affine torsors avoid that problem through
\BGtag{0245}; curve families require algebraic spaces.

\BGsolution{ex:change}
\hint{Track the order of right multiplications.}
The composite $R_{h_{jk}}\circ R_{h_{ij}}$ sends $x$ to
$xh_{ij}h_{jk}$, hence $h_{ik}=h_{ij}h_{jk}$. The new transition is
\[
 R_{b_j}\circ R_{h_{ij}}\circ R_{b_i}^{-1},
\]
which sends $x$ to $xb_i^{-1}h_{ij}b_j$. Therefore
\[
 h_{ij}'=b_i^{-1}h_{ij}b_j.
\]
On a triple overlap the adjacent factors $b_jb_j^{-1}$ cancel,
giving $h_{ij}'h_{jk}'=b_i^{-1}h_{ik}b_k=h_{ik}'$.
For $\bZ/2\bZ$ this becomes
$\varepsilon_{ij}'=b_i+\varepsilon_{ij}+b_j$; changing local labels
changes the transition functions but not the glued cover.

\BGsolution{ex:forget}
\hint{Use $\bC\otimes_{\bR}\bC\simeq\bC\times\bC$.}
The nonsplit cover is finite \emph{\'etale} of degree two, with conjugation
exchanging the two complex embeddings. Base changing it to $\bC$ gives
the algebra $\bC\otimes_{\bR}\bC\simeq\bC\times\bC$, so it becomes
the trivial torsor, equivariantly. Over $\bR$ it has no section: no
$\bR$-algebra map $\bC\to\bR$ exists. The split cover has sections,
so the covers are not isomorphic.

For (a), transport the canonical descent data of the nonsplit torsor
through its trivialization over $\bC$. This gives a datum of trivial
local torsors. Were it effective in the prestack of trivial torsors,
compatible uniqueness in the full torsor stack would make the nonsplit
torsor isomorphic to a split one, a contradiction. The arrow sheaves
of the trivial-torsor prestack still satisfy the sheaf axiom because
they are given by the sheaf $G$.

For (b), the two distinct global isomorphism classes restrict to the
same local class on this cover. Thus the isomorphism-class presheaf is
not even separated. The stack keeps the chosen isomorphisms and their
cocycles; passing just to classes discards precisely that information.

\BGsolution{ex:coproducts}
\hint{Disjoint union corresponds to a product of rings; fiber product
of affine schemes corresponds to a tensor product of rings.}
The three rings, in the stated order, are
\[
k[t]\times k[t],\qquad k[t_1,t_2],\qquad k[t].
\]
The last calculation uses $k[t]\otimes_{k[t]}k[t]\simeq k[t]$.
The fold corresponds to the diagonal homomorphism
$k[t]\to k[t]\times k[t]$, $f\mapsto(f,f)$.
Its composite with either projection of rings is the identity, so
the two component inclusions are sections of the fold. All maps
respect the map to $\Spec k$, giving the arrow in $\Sch/B$.
In $\Sch/S$ its target is $(S,\id_S)$ and its source is the
disjoint union equipped with the fold as structure map.

\BGsolution{ex:universalmap}
\hint{Existence of a comparison map is not universality of its source.}
The fiber product is $\bA^2_k$. The induced map is its origin,
corresponding to $k[t,x]\to k$ with $t,x\mapsto0$. This ring map
is not an isomorphism, so neither is the scheme map. If $W$, with its
two projections, satisfied the same universal mapping property for
every competing scheme, the comparison would be an isomorphism by
the two-factorization argument in Section~\ref{sec:categories}.
For the second part the pullback algebra is
\[
k[x]/(x^2-1)\simeq k\times k,
\]
because $(x-1)$ and $(x+1)$ are comaximal when $2$ is invertible.
Thus the fiber over $t=1$ is two copies of $\Spec k$.

\BGsolution{ex:testpoints}
\hint{A map into the affine line is determined by the image of its
coordinate, including nilpotent information.}
The ring maps $k[x]\to k[\epsilon]/(\epsilon^2)$ with images $0$
and $\epsilon$ are different, since $\epsilon\ne0$. Both induce
the same map on spectra: the unique prime $(\epsilon)$ contracts to
$(x)$. Thus topological-point data alone do not determine a morphism.
For any $k$-algebra map $\alpha:R\to R'$, one has
$\alpha(r^2)=\alpha(r)^2$, which is the naturality equation.
The identity point of $\bA^1_k$ corresponds to $x\in k[x]$.
Applying the rule gives $x^2$, so the recovered morphism is the
squaring map, defined by $k[x]\to k[x]$, $x\mapsto x^2$.

\BGsolution{ex:groupobjects}
\hint{The labels form a group, the copies of $B$ form a scheme, and
sections choose labels locally constantly.}
$H$ is the two-element abstract group. The scheme $H_B$ is two
copies $B_0\amalg B_1$ of the affine line, with the fold map to $B$.
The unit includes $B_0$, multiplication sends the component $(a,b)$
to $B_{a+b}$ by the identity on the base, and inversion is the
identity on each component. These are morphisms of $B$-schemes.
Since $B$ is nonempty and connected, $H_B(B)\simeq H$.
The two connected components of $T=B\amalg B$ can choose their
labels independently, giving $H_B(T)\simeq H\times H$.
Finally $H_B\times_B S\simeq S\amalg S$, with exactly the same
componentwise group operations over $S$.

\BGsolution{ex:fpflat}
\hint{Count algebra generators separately from module generators;
test flatness by tensoring an injection.}
The first algebra has one generator and one relation. Evaluation at
$0$ and $1$ identifies it with $A\times A$, which is free of rank
two as an $A$-module and hence flat. The algebra $A[x]$ has one
generator and no relations; its module basis is the infinite family
$1,x,x^2,\ldots$. Finitely many polynomials have bounded degrees,
so their $A$-linear span cannot contain all powers of $x$ when
$A\ne0$.
The algebra $\bZ/2\bZ$ is $\bZ/(2)$, hence finitely presented with
no variables and one relation. Tensor the injection
$\bZ\xrightarrow{\times2}\bZ$ with $\bZ/2\bZ$. The resulting
map is the zero endomorphism of a nonzero group and is not injective.
Thus finite presentation does not imply flatness.

\BGsolution{ex:coveroperations}
\hint{Ask what the new target is in each operation.}
The pulled-back cover of $T$ is
$\{T\times_S S_i\to T\}_i$; this is part~(3) of \BGtag{021O}.
The composite cover of $S$ is
$\{U_{ij}\to S_i\to S\}_{i,j}$; this is part~(2).
Tag~\texttt{021N} instead asserts inclusions between classes of
covers, so it does not supply these conclusions.
The pullback of the fold is $T\amalg T\to T$. Covering each of
these two sheets by two copies gives the four-copy fold onto $T$
after composition. Doing this instead to the original two $S$-sheets
gives the four-copy fold onto $S$. These are still covers, though a
four-copy fold over a nonempty base is not a double-cover torsor for
the two-element group. Being a cover and being that particular kind of torsor
are different conditions.
Finally the arrows of the site are all $B$-scheme morphisms.
For $B=\Spec\bZ$, the map $\Spec\bF_2\to B$ is an arrow but is
not flat, as Exercise~\ref{ex:fpflat} shows; it is not an fppf
covering map.

\begin{thebibliography}{9}
\raggedright
\bibitem{Stacks}
The Stacks Project Authors, \emph{The Stacks Project},
\href{https://stacks.math.columbia.edu}{stacks.math.columbia.edu}.
The linked tags in the text are permanent identifiers; statements and
numbering were checked in September 2026. Principal groups of references:
\begin{itemize}
\item Foundations: fiber products \BGtag{01JO}, \BGtag{01JQ};
product-ring spectra \BGtag{00ED}; relative affine space \BGtag{01M0},
\BGtag{01LX}; affine morphism--ring map correspondence \BGtag{01I1},
\BGtag{01I2}; functors of points \BGtag{001O}, \BGtag{01J5};
Yoneda \BGtag{001P};
sites \BGtag{00VH}; affine morphisms \BGtag{01S5}; local finite
presentation \BGtag{01TP}; flatness \BGtag{01U2}.
\item Topologies on Schemes, \S34.7 and \S34.9:
\BGtag{021M}, \BGtag{021N}, \BGtag{021O}, \BGtag{022B}, \BGtag{022C}.
\item Group schemes and actions: \BGtag{022S}, \BGtag{022T},
\BGtag{047F}, \BGtag{022Z}.
\item Groupoid Schemes, \S39.11, and Examples of Stacks, \S95.14--95.15:
\BGtag{049A}, \BGtag{0499}, \BGtag{049B}, \BGtag{049C},
\BGtag{04UR}, \BGtag{0CQJ}.
\item Stacks, \S8.2--8.5: \BGtag{026A}, \BGtag{02ZB},
\BGtag{026B}, \BGtag{026C}, \BGtag{02ZI}; quotient-stack definition:
\BGtag{044Q}.
\item Categories, \S4.33: \BGtag{02XK}, \BGtag{02XM}, \BGtag{02XO}.
\item Descent, \S35.13, \S35.23, \S35.35, \S35.37:
\BGtag{023Q}, \BGtag{02L4}, \BGtag{02L5}, \BGtag{02L2},
\BGtag{02KY}, \BGtag{02W0}, \BGtag{040L}, \BGtag{0245}.
\item Smooth and finite \emph{\'etale} examples:
\BGtag{02VL}, \BGtag{055U}, \BGtag{02VO}, \BGtag{02VN},
\BGtag{04HN}, \BGtag{00UB}, \BGtag{00UC}.
\item Algebraic spaces and curves: \BGtag{04UA}, \BGtag{0ADV},
\BGtag{041V}, \BGtag{041W}, \BGtag{0422}, \BGtag{04NS},
\BGtag{03MO}, \BGtag{04WP}, \BGtag{03XR}, \BGtag{06LR},
\BGtag{0D4Z}, \BGtag{0D1G}, \BGtag{0D52}.
\item Limits of scheme descent and nonscheme curve families:
\BGtag{08KF}, \BGtag{0D5E}.
\end{itemize}

\bibitem{Fantechi}
Barbara Fantechi,
\href{https://doi.org/10.1007/978-3-0348-8268-2_20}{\emph{Stacks for Everybody}}, in
\emph{European Congress of Mathematics}, Barcelona, 2000, Vol.~I,
Progress in Mathematics~201, Birkh\"auser, 2001, pp.~349--359.
Sections~3--4 develop categories
fibered in groupoids and the two stack conditions; Section~6.3 discusses
canonical pullback identifications. Local transcription:
\texttt{StacksForEverybody.tex} in the notebook's
\texttt{Papers/Stacks} directory.

\bibitem{Poonen}
Bjorn Poonen, \emph{Rational Points on Varieties}, Graduate Studies in
Mathematics~186, American Mathematical Society, 2017.
\href{https://math.mit.edu/~poonen/papers/Qpoints.pdf}{Author-hosted text},
Section~2.3, especially Definitions~2.3.1 and~2.3.3, for $T$-valued
points and the functor of points, and Section~2.3.1 for rational points.
\end{thebibliography}

\end{multicols}
\end{document}
